% =====================================================================
%  Chapter 26 -- Virtual Trees (Auxiliary Trees)                  WRITTEN
%  (folder ch30_virtual_trees; label prefix c30; PRINTS as
%  Chapter 26; folder number 30; build with:  python tools/build.py --only 30)
%
%  NUMBERING TRAP (the Ch.22/23/24/25 trap, mirrored): \chref{c30} prints
%  "Chapter 26" (THIS chapter).  Do NOT confuse the label with the printed
%  number: \chref{c23} prints "Chapter 25" (link-cut, the PREVIOUS chapter),
%  \chref{c21} prints "Chapter 27" (Mo, the NEXT chapter), \chref{c11} prints
%  "Chapter 16" (LCA -- the engine this chapter rides on), \chref{c24} prints
%  "Chapter 5" (Euler tour / tin-tout), \chref{c16} prints "Chapter 22" (HLD),
%  \chref{c18} prints "Chapter 24" (centroid), \chref{c10} prints "Chapter 13"
%  (tree DP), \chref{c09} prints "Chapter 12" (tree metrics), \chref{c28} prints
%  "Chapter 17" (DSU on tree), \chref{c20} prints "Chapter 15" (rerooting).
%  Every cite of THIS chapter is \chref{c30}.
%
%  Owns VIRTUAL / AUXILIARY TREES -- compressing a FIXED tree down to the
%  small set of nodes a query actually touches (the marked set plus the LCAs of
%  adjacent marked nodes in DFS order), building that compressed tree in O(k log k)
%  with a stack, bounding its size by 2k-1, then running a DP on it whose edges
%  carry the SKIPPED path lengths (dist = depth difference) and whose subtrees
%  carry the SKIPPED-node counts.  It STARTS where Ch.16 (LCA, \chref{c11}) ENDS:
%  ch11/99_chapter_end explicitly promises "Ch.26 (virtual trees) compresses a
%  k-node set via sort-by-tin + LCA of adjacent pairs; this chapter's
%  \thmref{c11-distf} and lifting are the inner loop."  THIS chapter discharges
%  that promise -- \lca, \dist, \tin and binary lifting are used, never re-derived.
%  Sections:
%    * 26.1 construction -- given k marked nodes, sort them by \tin{}, insert the
%      LCAs of adjacent pairs, dedup, and link each node to its nearest ancestor
%      in the set with a monotone stack; the \term{virtual tree} has at most
%      2k-1 nodes and every marked-node-to-marked-node path of the real tree is
%      preserved exactly (thm:c30-size).  Solved CF 2127E "Ancient Tree"
%      (2100, L5) -- build the virtual tree of EACH colour class, a vertex is a
%      forced "cutie" iff it lies in 2+ colour virtual trees; answer = sum of
%      weights of those vertices, reconstructed by a colouring DFS.  Live sample
%      hand-replayed (star test -> 0).
%    * 26.2 DP on the compressed tree -- an edge (u,v) with u an ancestor of v
%      carries weight \dist(u,v)=\depth(v)-\depth(u), so a tree DP over the
%      virtual tree reports the exact value a DP over the real tree would, while
%      touching only O(k) nodes; the nodes "skipped" inside a virtual edge still
%      have to be charged, and the per-subtree count of real nodes is carried
%      alongside.  Solved CF 2252F "Spectral Components" (2600, L6) -- per colour
%      build the virtual tree, run a bottom-up DP for m_in on each virtual
%      subtree, weight a virtual edge by the monotone centroid argument, then
%      greedily take the k_c-1 heaviest edges.  Live samples captured (S1->1,3,-1,-1,-1).
%    * 26.3 applications -- the two workhorse families: pairwise distance sums
%      over a marked set (each virtual edge contributes w * cnt * (k-cnt)) and
%      Steiner-like connected-subgraph counting on the marked set; both run in
%      O(k log k) after one LCA preprocess.  WORKED examples only, no solvedproblem.
%    * 26.4 recognition -- reading a statement and spotting the "small marked
%      set, big tree" signature that virtual trees compress; the constraint that
%      gives it away is sum k << n across queries, or a per-colour / per-subset
%      re-computation over the same fixed tree.  WORKED reading only.
%    * 26.5 the method boundary -- when a virtual tree is the WRONG tool: only a
%      FEW nodes are marked but the answer needs EVERY node's value (tree DP /
%      reroot); the terminals can be joined through NEW branch nodes (bitmask
%      Steiner DP, not a virtual tree); the whole tree is marked (virtual tree =
%      the tree, no gain); queries are on a tree that CHANGES (link-cut Ch.25).
%      WORKED examples only -- no solvedproblem (in-section count stays 2).
%  Chapter-end Summit: CF 1292D "Chaotic V." (2700, L6) -- the top of the band.
%  An infinite "divisor tree" where node x>1 attaches to x / (lowest prime
%  divisor of x); n mind-fragments sit at nodes k_i! (factorials).  Choose an
%  assembly node P minimising the sum of path lengths from every fragment to P
%  (a fragment repeated at a node counts multiple times).  The factorials
%  1!,2!,...,K! are already in DFS order of the divisor tree, so the marked set
%  is closed under \tin{} sorting; build the virtual tree over the O(n) distinct
%  factorials (at most 2n nodes), push each fragment's multiplicity onto its
%  factorial node, and the answer is the WEIGHTED CENTROID (1-median) sum on that
%  virtual tree -- edges carry dist(a!,b!) = Omega(a!)-Omega(b!) (prime-factor
%  counts).  All four official samples hand-replayed (5, 6, 6, 11).
%
%  Diagram-first: this chapter is TREE based.  The running example is the shared
%  rooted tree T used across sections (defined once, authoritatively, in
%  00_opener).  Rooted at 1, edges (parent -> child, 1-based):
%    1->2, 1->3, 2->4, 2->5, 3->6, 3->7, 5->8, 5->9, 7->10
%  (10 nodes).  The recurring marked set is S={4,8,6,10} (k=4); its virtual tree
%  is {1,2,3}+S with 6 compressed edges -- shown built, by stack, in 26.1.  The
%  Summit uses its OWN tiny divisor-tree cases (the CF samples), never T, so no
%  trace collides with the running example.
%
%  Anti-repetition contract (teach once, cross-reference by label):
%    * Ch.16 LCA (\chref{c11}) owns \lca, \dist (\thmref{c11-distf}), binary
%      lifting and the tin/tout ancestor test; 26.1-26.2 USE them as the inner
%      loop and NEVER re-derive them -- they are cited by \thmref/\secref.
%    * Ch.5 Euler tour (\chref{c24}) owns the \tin/\tout flatten and the
%      "u ancestor of v iff tin[u]<=tin[v]<tout[u]" test; 26.1 cites it for the
%      sort key, never re-teaches it.
%    * Ch.13 tree DP (\chref{c10}) owns the bottom-up subtree fold; 26.2 runs the
%      SAME fold on the compressed tree and says so once.
%    * Ch.25 link-cut (\chref{c23}) owns DYNAMIC forests; 26.5 draws the line:
%      virtual trees assume a FIXED tree, a moving topology is Ch.25's job.
%    * Explicitly OUT (prose + cite only): bitmask Steiner DP (26.5 boundary,
%      and CF 152E in the mixed set is the honest contrast); centroid distance
%      distributions (\chref{c18}); Mo on trees (\chref{c21}, the next chapter).
%
%  Known URL collisions (confirmed by grep of chapters/ + the live CF API;
%  prose-cite only, NEVER re-link -- urlsweep must print an empty DUPLICATES
%  section):
%    * CF 613D "Kingdom and its Cities" (2800) -> Ch.27 (\chref{c21_tree_dp_ii},
%      folder ch27).  The most famous virtual-tree problem; its URL lives there,
%      PROSE only.  Its rating (2800) is also above this chapter's band.
%    * CF 809E "Surprise me!" (3100) -> Ch.25 (\chref{c23}) mixed set, framed
%      THERE as "this is Chapter 26's territory".  Not re-linked here.
%    * CF 1681F "Unique Occurrences" (2300) -> Ch.17 DSU-on-tree (\chref{c28}).
%    * CF 1923E "Count Paths" (2000) and CF 915F (2400) -> Ch.20/Ch.27.  The USACO
%      virtual-tree list also names these; all live elsewhere, PROSE only.
%    * CF 1338B (Ch.13), CF 342E (Ch.12) -- never touched here.
%    * The three roster URLs below (2127E, 2252F, 1292D) and every mixed-set URL
%      were confirmed ABSENT from every other chapter by the grep sweep.
%
%  Roster (every title/rating/tag/sample live-verified against the Codeforces
%  problemset API / problem page / official editorial BEFORE a trace was written;
%  CF-forward):
%    solved: CF 2127E Ancient Tree        (2100, L5) (26.1) -- per-colour virtual
%            tree + membership count; live 4-case sample, star test replayed -> 0.
%            CF 2252F Spectral Components (2600, L6) (26.2) -- per-colour VT +
%            bottom-up DP + greedy; live samples -> 1 3 -1 -1 -1 / 3 0 -1.. / 3 2..
%    worked: the stack build of the VT of S={4,8,6,10} on T (26.1); the
%            edge-weight/skipped-count carry (26.2); the pairwise distance-sum
%            fold w*cnt*(k-cnt) (26.3); a connected-subgraph count on the VT (26.3);
%            the "small marked set" recognition reads (26.4); the six-tool pick (26.5).
%    summit: CF 1292D Chaotic V.          (2700, L6) -- divisor tree, fragments at
%            k_i!, minimise sum of distances to assembly node P; factorials are DFS
%            order, build the VT over distinct factorials, weighted centroid.  All
%            four samples hand-replayed (5, 6, 6, 11).
%    mixed (10, ascending by live-verified rating, every ID confirmed ABSENT from
%            every other chapter via the grep sweep -- locked at the
%            99_chapter_end write): 1901E(2200), 2253E(2200), 2112E(2200),
%            2254G(2200), 2247E(2300), 1494D(2300), 2194F1(2300), 1633E(2400),
%            152E(2500), 1254D(2700).  A spread of genuine VT-representation
%            drills and deliberate REJECTION cases -- their hints name the TRUE
%            tool (tree DP, bitmask Steiner DP, distance-matrix reconstruction,
%            binary search + DSU) that beats a virtual tree.
% =====================================================================
\bookchapter{Virtual Trees (Auxiliary Trees)}{c30}
\chapterinfo{CF 2200--2700}{Written}

% =====================================================================
%  Chapter 26 opener -- "Chapter at a Glance", three reading paths + the
%  ladder staircase, THE running example (the shared 10-node rooted tree T
%  plus the recurring marked set S={4,8,6,10} and its virtual tree), the
%  who-owns-what comparison against Ch.16 (LCA), Ch.5 (Euler/tin), Ch.13
%  (tree DP), Ch.25 (link-cut) and Ch.22 (HLD), the recognition guide, a
%  constraint->technique table, and notation.
%  Exports fig:c30-tree, fig:c30-taxonomy.
%  No \term{lca} / \term{tin} here -- they are indexed by Ch.16/Ch.5; \term the
%  NEW objects "virtual tree", "auxiliary tree", "marked set", "compressed
%  edge" only.
% =====================================================================

\begin{chaptertheorybox}{Chapter at a Glance}
Chapter~16 (\chref{c11}) gave you \(\lca\): one DFS fixes \(\tin\), \(\tout\),
\(\depth\), and a binary-lifting table turns every ``lowest common ancestor'' and
``distance between two nodes'' into \(O(\log n)\) arithmetic
(\(\thmref{c11-distf}\)).  It closed by pointing forward -- ``a \emph{virtual tree}
(Ch.~26) compresses a \(k\)-node set via sort-by-tin plus the LCA of adjacent
pairs; this chapter's distance formula and lifting are the inner loop.''  This
chapter cashes that cheque.

A virtual tree answers a question no full tree DP can afford: you have one
\emph{fixed} tree on \(n\) nodes and a \emph{small set of \(k\) marked nodes}
(a colour class, a set of query terminals, a set of special cities), and you must
run some DP \emph{over just those \(k\) nodes and the paths between them} -- once
per colour, per query, per subset.  Doing it on the real tree is \(O(n)\) each
time; the virtual tree collapses the whole structure to at most \(2k-1\) nodes in
\(O(k\log k)\) -- the marked nodes plus the LCAs that connect them -- while an edge
keeps the \emph{length} of the path it swallowed, so every distance, subtree-size
and connectivity fact among the marked nodes survives untouched.

\textbf{You will learn to}
\begin{itemize}[nosep]
  \item build the \term{virtual tree} of a \term{marked set}: sort the \(k\) nodes by
        \(\tin\), insert \(\lca\) of each adjacent pair, dedup, then link every node
        to its nearest kept ancestor with a monotone stack -- \(\bigO{k\log k}\) and
        at most \(2k-1\) nodes (\(\secref{c30-construction}\), \thmref{c30-size},
        solved CF 2127E \emph{Ancient Tree}); \hfill\lv{3}--\lv{4}
  \item run a DP on the compressed tree, charging each \term{compressed edge} its real
        length \(\dist(u,v)=\depth(v)-\depth(u)\) and carrying the count of swallowed
        nodes, so the answer equals the \(O(n)\) DP done honestly (\(\secref{c30-dp}\),
        solved CF 2252F \emph{Spectral Components}); \hfill\lv{5}--\lv{6}
  \item spend the machine on the two workhorse families -- the sum of pairwise
        distances over a marked set (each edge contributes \(w\cdot\mathit{cnt}\cdot
        (k-\mathit{cnt})\)) and Steiner-like connected-subgraph counts
        (\(\secref{c30-apps}\)); \hfill\lv{5}
  \item read a statement and recognise the ``big tree, tiny marked set, re-run per
        colour/subset'' signature (\(\secref{c30-recognition}\)); \hfill\lv{4}
  \item know when a virtual tree is the WRONG tool -- bitmask Steiner DP, whole-tree
        DP, or a topology that changes (\(\secref{c30-boundary}\), the pick). \hfill\lv{5}
\end{itemize}

\textbf{Prerequisites.} Chapter~16 (\chref{c11}) -- \(\lca\), \(\dist\), binary
lifting: the inner loop, used and never re-derived.  Chapter~5 (\chref{c24}) -- the
\(\tin/\tout\) Euler order that is the sort key.  Chapter~13 (\chref{c10}) -- the
bottom-up subtree fold you re-run on the compressed tree.  Chapter~12 (\chref{c09})
-- \(\depth\) as the currency a compressed edge spends.  Chapter~25 (\chref{c23})
and Chapter~22 (\chref{c16}) -- the neighbours on either side: a link-cut tree is
this idea with a \emph{moving} topology; HLD is a fixed decomposition that does not
shrink with \(k\).
\end{chaptertheorybox}

% ---------------- Reading paths + ladder position ----------------
\begin{chaptertheorybox}{How to Read This Chapter (three paths)}
Level badges \(\lv3\)--\(\lv6\): \(\lv3\) is hand-building the virtual tree of a
small marked set and reading off its \(2k-1\) size; \(\lv{4}\) is the monotone-stack
construction coded and the \(\bigO{k\log k}\) bound justified; \(\lv{5}\) is a DP
on the compressed tree that charges compressed-edge lengths and swallowed-node
counts, or the pairwise-distance fold; \(\lv{6}\) is the Summit -- a virtual tree
over factorial nodes in an infinite divisor tree, solved by a weighted centroid.

\textbf{Construction path (\(\lv3\)--\(\lv{5}\), ``make the tree shrink'').}
\secref{c30-construction} in full -- sort-by-tin, adjacent LCAs, the stack,
\thmref{c30-size}, solved CF 2127E -- then \secref{c30-dp}, where the compressed
tree is fed a DP (solved CF 2252F), capped by the Summit \(\probref{c30-summit}\).

\textbf{Full-power path (\(\lv{4}\)--\(\lv{6}\), ``a DP over a \(k\)-set on a
\(10^5\) tree'').} \secref{c30-dp} for edge weights and the skipped-count carry,
\secref{c30-apps} for the distance-sum and connected-count families, then the
Summit, where the marked set is the factorial nodes and the payoff is a weighted
1-median.

\textbf{Boundary path (\(\lv{4}\)--\(\lv{5}\), ``is this even a virtual-tree job?'').}
\secref{c30-recognition} and \secref{c30-boundary}: the tell (tiny \(k\), re-run per
colour) and the four look-alikes that are not -- bitmask Steiner DP, whole-tree DP,
a moving forest for the link-cut tree (\chref{c23}), a static range for HLD
(\chref{c16}).
\end{chaptertheorybox}

\begin{figure}[H]
\centering
\begin{tikzpicture}[diagram, x=8.5mm, y=6.5mm]
  \foreach \st/\lab in {1/lv1,2/lv2,3/lv3,4/lv4,5/lv5,6/lv6} {
    \fill[\lab!22, draw=\lab, line width=0.8pt] (\st-1,0) rectangle (\st,\st*0.52);
    \node[font=\scriptsize\bfseries, text=\lab] at (\st-0.5,\st*0.52-0.42) {L\st};}
  \node[annot, anchor=south west, font=\scriptsize]
    at (0.05,0.62) {prereq (Ch.16): LCA, dist, tin -- the inner loop};
  \node[annot, anchor=south west, font=\scriptsize]
    at (1.05,1.14) {26.1 sort by tin + adjacent LCAs; dedup; the 2k-1 bound};
  \node[annot, anchor=south west, font=\scriptsize]
    at (2.05,1.66) {26.1 stack build, solved CF 2127E (per-colour VT)};
  \node[annot, anchor=south west, font=\scriptsize]
    at (3.05,2.18) {26.2 compressed-edge weight, skipped counts, solved CF 2252F};
  \node[annot, anchor=south west, font=\scriptsize, text=lv5]
    at (4.05,2.70) {26.3 distance-sum / connected-count; 26.4 recognition};
  \node[annot, anchor=south west, font=\scriptsize, text=lv6]
    at (5.05,3.22) {26.5 the pick; Summit CF 1292D (factorial VT + weighted centroid)};
\end{tikzpicture}
\caption{Chapter~26 on the six-step ladder: hand-building a virtual tree and
reading its size is \(\lv3\); the coded monotone-stack build and the \(O(k\log k)\)
bound are \(\lv{4}\); a DP that charges compressed-edge lengths and swallowed-node
counts (or the pairwise-distance fold) is \(\lv{5}\); the Summit (a virtual tree over
factorial nodes solved by a weighted centroid) is step~6.}
\end{figure}

% ---------------- The running example ----------------
\subsection*{The Running Example: One Fixed Tree, One Tiny Marked Set}

Every section in this chapter runs on the same rooted tree \(T\) on \(n=10\)
nodes, fixed once here and authoritative (if a later trace disagrees, this figure
wins).  It is a single tree rooted at \(1\); parent-to-child edges, 1-based:
\[
  1\!\to\!2,\quad 1\!\to\!3,\quad 2\!\to\!4,\quad 2\!\to\!5,\quad
  3\!\to\!6,\quad 3\!\to\!7,\quad 5\!\to\!8,\quad 5\!\to\!9,\quad 7\!\to\!10 .
\]
A left-to-right DFS assigns entry times \(\tin\): node \(1\!:\!1\), \(2\!:\!2\),
\(4\!:\!3\), \(5\!:\!4\), \(8\!:\!5\), \(9\!:\!6\), \(3\!:\!7\), \(6\!:\!8\),
\(7\!:\!9\), \(10\!:\!10\).  Depths are \(\depth(1)=0\);
\(\depth(2)=\depth(3)=1\); \(\depth(4)=\depth(5)=\depth(6)=\depth(7)=2\);
\(\depth(8)=\depth(9)=\depth(10)=3\).  The recurring \term{marked set} is
\(S=\{4,8,6,10\}\) (\(k=4\)) -- think of them as the four special vertices a query
cares about.  Their virtual tree \(A(S)\) keeps \(S\) plus the LCAs of adjacent
pairs in \(\tin\)-order, \(\{\lca(4,8),\lca(8,6),\lca(6,10)\}=\{2,1,3\}\), giving
\(7=2k-1\) nodes and six \term{compressed edges}; nodes \(5,7,9\) vanish into edge
lengths (\(\secref{c30-construction}\)).

\begin{figure}[H]
\centering
\begin{tikzpicture}[diagram, scale=0.9, transform shape,
   tnode/.style={circle, draw=primaryblue!70, fill=primaryblue!8, line width=0.7pt,
                 minimum size=6.2mm, font=\scriptsize\bfseries, inner sep=1pt},
   mk/.style={circle, draw=warningred, fill=warningred!16, line width=1.1pt,
             minimum size=6.8mm, font=\footnotesize\bfseries, inner sep=1pt},
   kee/.style={circle, draw=successgreen!75, fill=successgreen!12, line width=0.9pt,
             minimum size=6.2mm, font=\scriptsize\bfseries, inner sep=1pt},
   drop/.style={circle, draw=gray!60, fill=gray!8, line width=0.5pt, dashed,
             minimum size=5.4mm, font=\scriptsize\footnotesize, inner sep=1pt, text=gray}]
  % left: the real tree T
  \node[mk] (t1) at (0,0) {1};
  \node[kee] (t2) at (-2.2,-1.3) {2};
  \node[kee] (t3) at (2.2,-1.3) {3};
  \node[mk] (t4) at (-3.4,-2.7) {4};
  \node[drop] (t5) at (-1.2,-2.7) {5};
  \node[mk] (t6) at (1.5,-2.7) {6};
  \node[drop] (t7) at (3.3,-2.7) {7};
  \node[mk] (t8) at (-2.0,-4.1) {8};
  \node[drop] (t9) at (-0.5,-4.1) {9};
  \node[mk] (t10) at (3.3,-4.1) {10};
  \draw (t1)--(t2) (t1)--(t3) (t2)--(t4) (t2)--(t5) (t3)--(t6) (t3)--(t7)
             (t5)--(t8) (t5)--(t9) (t7)--(t10);
  \node[annot, font=\scriptsize, anchor=south] at (0.4,0.35) {the fixed tree T (n=10)};
  % right: the virtual tree
  \node[kee] (v1) at (8.2,0) {1};
  \node[kee] (v2) at (6.6,-1.6) {2};
  \node[kee] (v3) at (9.8,-1.6) {3};
  \node[mk] (v4) at (5.6,-3.3) {4};
  \node[mk] (v8) at (7.4,-3.3) {8};
  \node[mk] (v6) at (9.0,-3.3) {6};
  \node[mk] (v10) at (10.8,-3.3) {10};
  \draw[very thick] (v1)--node[left,font=\scriptsize,text=gray]{1} (v2)
        (v1)--node[right,font=\scriptsize,text=gray]{1} (v3)
        (v2)--node[left,font=\scriptsize,text=gray]{1} (v4)
        (v2)--node[right,font=\scriptsize,text=gray]{2} (v8)
        (v3)--node[left,font=\scriptsize,text=gray]{1} (v6)
        (v3)--node[right,font=\scriptsize,text=gray]{2} (v10);
  \node[annot, font=\scriptsize, anchor=south] at (8.4,0.35) {virtual tree A(S), 2k-1=7 nodes};
\end{tikzpicture}
\caption{The running example.  Left: the fixed tree \(T\); red nodes are the marked
set \(S=\{4,8,6,10\}\), green nodes are the LCAs the construction must add
(\(\{1,2,3\}\)), grey dashed nodes \(\{5,7,9\}\) are \emph{swallowed}.  Right: the
virtual tree \(A(S)\) -- every real edge between two kept nodes is replaced by a
\term{compressed edge} (thick) carrying the number of real edges it swallowed
(\(2\!\to\!8\) has length \(2\) along \(2\text{-}5\text{-}8\)).  The distance and
ancestor relation between any two red/green nodes is exactly the real tree's.}
\label{fig:c30-tree}
\end{figure}

% ---------------- Which job, cost chart ----------------
\subsection*{Which Job? (the question-shape chart)}

\begin{center}
\small
\renewcommand{\arraystretch}{1.25}
\setlength{\tabcolsep}{3pt}
\begin{tabular}{@{}p{0.31\linewidth} p{0.17\linewidth} p{0.28\linewidth} >{\raggedright\arraybackslash}p{0.14\linewidth}@{}}
\toprule
\textbf{Question} & \textbf{Technique} & \textbf{Why not a full-tree DP} & \textbf{Where} \\
\midrule
a DP over \(k\) marked nodes, re-run per colour/query & virtual tree + DP & \(O(n)\) per copy blows the budget; VT is \(O(k\log k)\) & \secref{c30-construction} \\
sum of \(\dist\) over all pairs of a \(k\)-set & VT, edge \(\cdot\)cnt\(\cdot(k\)-cnt) & pair walk is \(O(kn)\); VT touches \(2k\) edges & \secref{c30-apps} \\
cheapest connected subgraph spanning \(k\) terminals & VT + fold on it & the terminals' Steiner tree has \(O(k)\) branching nodes & \secref{c30-dp} \\
for a set of nodes, is \(u\) an ancestor of \(v\)? & \(\tin/\tout\) test & free once \(\tin\) is known (Ch.5) & \secref{c30-construction} \\
terminals joined through \emph{new} branch nodes & \emph{bitmask Steiner DP} & VT only keeps existing nodes, cannot invent terminals & \secref{c30-boundary} \\
the tree \emph{changes} between queries & \emph{go to link-cut} & VT assumes a fixed tree; edits need Ch.25 & \chref{c23} \\
\bottomrule
\end{tabular}
\end{center}

\begin{insightbox}[One promise decides everything: keep the nodes that branch]
A virtual tree keeps exactly the marked nodes and the nodes where the paths
between them \emph{branch} -- and a branch point of marked nodes is always an LCA of
two of them.  Sorting the \(k\) marked nodes by \(\tin\) makes every needed LCA the
LCA of two \emph{adjacent} marked nodes, so \(k-1\) LCA calls (plus dedup) already
close the set under ancestry.  Compress every maximal chain of unmarked, non-branch
nodes into a single weighted edge, and you get a tree with at most \(2k-1\) nodes
that answers every distance, ancestor and connectivity question among the marked
set \emph{exactly} -- \thmref{c30-size}.  The whole chapter is one idea: \emph{make
the tree shrink to \(O(k)\) without changing what the \(k\) nodes know about each
other.}
\end{insightbox}

% ---------------- Taxonomy chart ----------------
\begin{figure}[H]
\centering
\begin{tikzpicture}[diagram, scale=0.95, transform shape,
    mnode/.style={circle, draw=#1, line width=0.9pt, fill=#1!8, minimum size=9mm,
                  font=\scriptsize\bfseries, text=#1!60!black}]
  \node[mnode=primaryblue] (a) at (0,0)   {close};
  \node[annot, font=\scriptsize, anchor=west, text width=9.8cm, align=left] at (1.7,0) {sort the k marked nodes by tin; for each adjacent pair take lca and add it; dedup. This set A(S) is closed under lca, and has at most 2k-1 nodes (26.1)};
  \node[mnode=successgreen] (b) at (0,-1.1) {link};
  \node[annot, font=\scriptsize, anchor=west, text width=9.8cm, align=left] at (1.7,-1.1) {sweep A(S) by tin with a monotone stack: pop until the stack top is an ancestor of the current node, then make the current node a virtual child of the top. Edge weight = depth difference (26.1--26.2)};
  \node[mnode=difforange]   (c) at (0,-2.2) {fold};
  \node[annot, font=\scriptsize, anchor=west, text width=9.8cm, align=left] at (1.7,-2.2) {run the usual bottom-up tree DP on the compressed tree, but charge each virtual edge its real length and carry the count of swallowed nodes, so the value equals the honest O(n) DP (26.2, 26.3; solved CF 2252F)};
  \node[mnode=warningred]   (d) at (0,-3.3) {recognise};
  \node[annot, font=\scriptsize, anchor=west, text width=9.8cm, align=left] at (1.7,-3.3) {the tell is a fixed big tree, a small marked set, and a computation repeated per colour / per query / per subset with sum k << total budget (26.4)};
  \node[mnode=softpurple]   (e) at (0,-4.4) {reject};
  \node[annot, font=\scriptsize, anchor=west, text width=9.9cm, align=left] at (1.7,-4.4) {whole tree needed -> tree DP; terminals can meet at NEW nodes -> bitmask Steiner DP; topology moves -> link-cut (Ch.25); static ranges -> HLD/Euler (Ch.22/Ch.5) (26.5)};
\end{tikzpicture}
\caption{The chapter's five-state taxonomy.  \emph{close} turns the marked set into
an LCA-closed set; \emph{link} builds the compressed tree with a stack; \emph{fold}
runs a weighted DP on it; \emph{recognise} spots the signature in a statement;
\emph{reject} hands the job to the right neighbour.  Each solved problem sits on one
row: CF 2127E on \emph{close}+a count, CF 2252F on \emph{fold}, the Summit CF 1292D
on \emph{link}+a weighted centroid.}
\label{fig:c30-taxonomy}
\end{figure}

% ---------------- Who owns what ----------------
\subsection*{Who Owns What (so nothing is learned twice)}

\begin{center}
\small
\renewcommand{\arraystretch}{1.25}
\setlength{\tabcolsep}{3pt}
\begin{tabular}{@{}p{0.28\linewidth} p{0.24\linewidth} p{0.44\linewidth}@{}}
\toprule
\textbf{Idea} & \textbf{Owner} & \textbf{This chapter's share} \\
\midrule
\(\lca\), \(\dist\), binary lifting & Ch.~16 (\chref{c11}) & the inner loop of \emph{close} and \emph{link}; used, cited by \thmref{c11-distf}, never re-derived \\
\(\tin/\tout\) flatten, ancestor test & Ch.~5 (\chref{c24}) & the sort key and the stack's ``is ancestor'' predicate; never re-taught \\
bottom-up subtree fold & Ch.~13 (\chref{c10}) & the same fold, run on the \(O(k)\) compressed tree with edge weights (26.2) \\
\(\depth\) as a distance currency & Ch.~12 (\chref{c09}) & a compressed edge's weight is a \(\depth\) difference (26.1--26.2) \\
dynamic forest, moving topology & Ch.~25 (\chref{c23}) & explicitly out; 26.5 routes ``edges change'' to the link-cut tree \\
static path/subtree \emph{ranges} & Ch.~22 HLD / Ch.~5 Euler & explicitly out; those do not shrink with \(k\) and are the boundary's first rejects \\
marked-set queries via lifting & Ch.~16 Summit & 26.x reuses that set-LCA insight but builds the whole induced tree, not one range \\
\bottomrule
\end{tabular}
\end{center}

% ---------------- Recognition guide ----------------
\begin{recognitionbox}[Pattern card -- ``a big fixed tree, a tiny set I re-run on'']
\reckey{statements that fix one tree of \(n\le 2\cdot10^5\) nodes and then hand you
a \emph{small} set of special vertices -- one colour class, one query's terminal list,
one subset -- and ask you to compute something over exactly that set and the paths
linking it (a DP, a distance sum, a min connected subgraph, ``for each colour over its
vertices\ldots'').  The tell is that the same \(O(n)\) tree DP would be re-run
\emph{many times}, once per subset, and the total number of marked vertices across all
runs is small -- \(\sum k \ll n\cdot(\text{runs})\).}
\recuse{
  \item per-colour / per-subset DP over that subset's induced subtree -- virtual tree + fold (\(\secref{c30-dp}\));
  \item sum of distances over all pairs of a \(k\)-set -- virtual tree, per-edge \(w\cdot\mathit{cnt}\cdot(k-\mathit{cnt})\) (\(\secref{c30-apps}\));
  \item minimum connected subgraph spanning \(k\) marked nodes (Steiner \emph{on a tree}) -- virtual tree DP (\(\secref{c30-apps}\)).
}
\recnotuse{
  \item the answer needs every node's own value, not just the marked ones -- full tree DP / rerooting (\chref{c10}, \chref{c20});
  \item terminals may be joined at nodes that are \emph{not} in the tree you were given / new branch nodes -- bitmask Steiner DP, not a virtual tree (\(\secref{c30-boundary}\));
  \item the tree itself is edited between queries -- link-cut tree (\chref{c23});
  \item you query ranges on a \emph{fixed} tree and nothing shrinks with \(k\) -- HLD / Euler+Fenwick (\chref{c16}, \chref{c24}).
}
One-line test: \emph{one fixed tree, a marked set of size \(k\ll n\), a computation
repeated over that set -- then collapse to \(O(k)\) with a virtual tree; if every node
matters or the set never shrinks, keep the full-tree engine.}
\end{recognitionbox}

% ---------------- Constraint -> technique ----------------
\subsection*{Reading the Constraints (which technique the limits want)}

\begin{center}
\small
\renewcommand{\arraystretch}{1.25}
\begin{tabular}{@{}p{0.34\linewidth} p{0.40\linewidth} p{0.20\linewidth}@{}}
\toprule
\textbf{Constraint signature} & \textbf{What it rules in/out} & \textbf{Technique} \\
\midrule
fixed tree \(n\le2\!\cdot\!10^5\), and \(\sum(\text{marked})\) over all runs \(\le 2\!\cdot\!10^5\) & a full \(O(n)\) DP per run is too slow; each run wants \(O(k\log k)\) & \secref{c30-construction} \\
``for each colour / each query, over its \(k\) vertices\ldots'' & the object lives on the marked set's induced subtree, not the whole tree & \secref{c30-dp} \\
sum of distances over all pairs of a small set & \(O(k^2)\) pairs too slow; fold per virtual edge & \secref{c30-apps} \\
Steiner tree / min connected subgraph on a TREE (terminals are existing nodes) & the Steiner nodes are already the marked set + their LCAs & \secref{c30-apps} \\
terminals may meet at \emph{new} nodes, \(k\le18\) & a virtual tree cannot create terminals; need bitmask DP & \secref{c30-boundary} \\
\bottomrule
\end{tabular}
\end{center}

% ---------------- Notation ----------------
\subsection*{Notation Used Only Here}

The fixed tree is \(T=(V,E)\), \(|V|=n\), rooted at \(1\).  A \term{marked set}
\(S\subseteq V\), \(|S|=k\), is the small set of vertices a single run cares about.
The \term{virtual tree} \(A(S)\) is the minimal ancestor-closed set containing \(S\)
(equivalently \(S\) together with \(\lca(x,y)\) for adjacent \(x,y\) in \(\tin\)-order),
viewed as a tree whose edges are \term{compressed edges}: a virtual edge \((u,v)\),
\(u\) a proper ancestor of \(v\), stands for the whole real path \(u\!\to\!v\) and
carries weight \(\dist(u,v)=\depth(v)-\depth(u)\).  \(\mathit{cnt}(v)\) is the number
of marked nodes in \(v\)'s virtual subtree; \(\size(v)\) is the number of \emph{real}
tree nodes swallowed inside \(v\)'s virtual subtree (used when the DP must charge
skipped nodes).  \(\Omega(x)\) is the number of prime factors of \(x\) with
multiplicity -- it is the depth of node \(x\) in the Summit's divisor tree.
\(\depth,\tin,\tout,\lca,\dist\) are Ch.~5/Ch.~16's, used unchanged.  Every section
runs on the \(T\) and \(S=\{4,8,6,10\}\) above; if a table here and a trace there
disagree, this opener is authoritative.

% =====================================================================
%  26.1  Construction -- close a marked set under LCA and link it with a
%  monotone stack.  Exports: sec:c30-construction  def:c30-vt
%  thm:c30-size  fig:c30-build  lst:c30-build  prob:c30-ancient
%  lst:c30-ancient
%  Solved CF 2127E "Ancient Tree" (2100, L5), live-verified against the CF
%  problem page + the Round 1041 editorial (which explicitly builds, per colour,
%  A = S_c u {lca(u,v)} and counts how many colour virtual trees each vertex is a
%  branch of).  Official 4-test sample read live; test 1 (4-node star) replayed
%  here to cost 0.  DISCHARGES the Ch.16 (sec:c11-lifting) forward promise "a virtual tree
%  compresses a k-node set via sort-by-tin + LCA of adjacent pairs"; LCA and tin
%  are USED, never re-derived.
% =====================================================================
\section{Building a Virtual Tree}\label{sec:c30-construction}

\begin{previouslybox}[Ch.~16 handed you \(\lca\) and promised this chapter]
Chapter~16 (\(\secref{c11-lifting}\)) proved that one DFS fixes \(\tin,\tout,\depth\), that
\(u\) is an ancestor of \(v\) iff \(\tin(u)\le\tin(v)<\tout(u)\), and that binary
lifting answers \(\lca(u,v)\) -- hence \(\dist(u,v)=\depth(u)+\depth(v)-2\depth(\lca)\)
(\(\thmref{c11-distf}\)) -- in \(O(\log n)\).  It then promised, in its own words,
that ``a virtual tree (Ch.~26) compresses a \(k\)-node set via sort-by-tin plus the
LCA of adjacent pairs; this chapter's distance formula and lifting are the inner
loop.''  This section is that compression.  We never rebuild a lifting table and
never re-prove the distance formula; we call \(\lca\) \(O(k)\) times and spend the
result on a set of \(k\ll n\) marked vertices.
\end{previouslybox}

\begin{topicintro}
Given the fixed tree \(T\) and a marked set \(S\) of size \(k\), the virtual tree is
the smallest subtree-shaped \emph{view} of \(T\) that (i) contains every marked node,
(ii) is closed under \(\lca\) (so all the branching is preserved), and (iii) replaces
every maximal chain of unmarked, non-branching nodes by one weighted edge.  Building
it has exactly two steps.  \emph{Close}: sort \(S\) by \(\tin\), add \(\lca\) of each
adjacent pair, dedup.  \emph{Link}: sweep the closed set by \(\tin\) with a stack of
the current root-to-node path, hanging each node on its nearest kept ancestor.  The
miracle, proved in \thmref{c30-size}, is that \(2k-1\) nodes is a hard ceiling -- the
structure is linear in the marked set, not in \(n\).

\textbf{Prerequisites.} Chapter~16 (\(\secref{c11-paths}\)) -- \(\lca\), \(\dist\).
Chapter~5 (\(\secref{c24-times}\), \chref{c24}) -- the \(\tin/\tout\) order and the
ancestor test, used as the sort key and the stack predicate.  Nothing else.

\textbf{After this section you can}
\begin{itemize}[nosep]
  \item close a marked set under \(\lca\) with \(k-1\) adjacent-pair calls; \hfill\lv{3}
  \item prove the closed set has at most \(2k-1\) nodes; \hfill\lv{4}
  \item build the compressed tree with one monotone-stack sweep in \(\bigO{k\log k}\); \hfill\lv{4}
  \item recognise a per-colour / per-subset problem and collapse it -- solved CF 2127E. \hfill\lv{4}
\end{itemize}
\end{topicintro}

% ---------------------------------------------------------------
\subsection{Closing the marked set under \(\lca\)}
% ---------------------------------------------------------------

\begin{definition}{Virtual (auxiliary) tree of a marked set}{c30-vt}
Let \(S\subseteq V(T)\) be the marked set.  The \term{virtual tree} \(A(S)\) is the
set \(A(S)=S\cup\{\lca(x,y):x,y\in S\}\) made into a tree by drawing an edge
\((u,v)\) whenever \(u\) is a proper ancestor of \(v\) in \(T\) and no element of
\(A(S)\) lies strictly between them on the real path.  Equivalently, \(A(S)\) is the
minimal \(\lca\)-closed superset of \(S\), and it is exactly the union of all simple
paths of \(T\) between pairs of marked nodes, with maximal unbranched chains
compressed into single weighted edges.
\end{definition}

Naively \(|A(S)|\) looks like it could be as large as the whole tree, because there
are \(\binom{k}{2}\) pairs.  The load-bearing fact is that only \emph{adjacent}
pairs in \(\tin\)-order are needed to close the set.

\begin{theorem}{Adjacent-pair LCAs suffice, and \(|A(S)|\le 2k-1\)}{c30-size}
Sort \(S\) by entry time as \(x_1,x_2,\dots,x_k\).  Then
\[
   S\cup\{\lca(x_i,x_{i+1}): 1\le i<k\}
\]
is already closed under \(\lca\) -- it equals \(A(S)\).  In particular \(|A(S)|\le
2k-1\).
\end{theorem}
\begin{proofbox}
Take any two marked nodes \(a,b\) and let \(w=\lca(a,b)\).  In the DFS order the
marked nodes in \(w\)'s subtree form a contiguous block, and \(a,b\) being under
\(w\) in different child-subtrees of \(w\) (or one being \(w\)) means there is some
\emph{adjacent} pair \(x_i,x_{i+1}\) inside that block whose paths leave \(w\)
through two different children (or one is \(w\) itself); for that pair
\(\lca(x_i,x_{i+1})=w\).  Hence every \(\lca\) of marked nodes appears as an
adjacent-pair \(\lca\), so the \(k-1\) adjacent LCAs already contain all of
\(A(S)\setminus S\).  For the size: the compressed tree has \(k\) marked leaves
(every element of \(S\) that is not an LCA of others is a leaf; add back those that
are internal, at most \(k\)) and an \(\lca\)-closed tree in which every internal node
has at least two children drawn from the kept set, so the number of internal nodes
is at most \(k-1\); together \(|A(S)|\le k+(k-1)=2k-1\). \(\Box\)
\end{proofbox}

\begin{insightbox}[\(\tin\)-adjacency is doing the work of \(\binom{k}{2}\) pairs]
A branching node of the marked set is, by definition, an LCA where the marked nodes
below it split into two different DFS child-buckets.  Sorting by \(\tin\) lines those
buckets up consecutively, so the split is always visible between two \emph{neighbours}
in the order -- which is why \(k-1\) LCA calls, not \(\binom{k}{2}\), close the set.
The bound \(2k-1\) is tight: take \(k\) leaves all sharing one root; their LCAs collapse
to the single root, giving \(k+1\); take a caterpillar of nested pairs and you approach
\(2k-1\).
\end{insightbox}

% ---------------------------------------------------------------
\subsection{Linking with a monotone stack}
% ---------------------------------------------------------------

Once the closed set \(C=A(S)\) is known and sorted by \(\tin\), each node must be
hung on its nearest kept ancestor.  Walking up with \(\lca\) per node is \(O(k\log
n)\) and fiddly; the clean way is one stack sweep that maintains the current
root-to-node path inside \(C\).

\begin{lstlisting}[caption={Build the virtual tree of a marked set. Input: marked nodes, and lca/tin from a preprocessing pass (Ch.16). Sort by tin, close under adjacent-pair lca, dedup, then a monotone stack links each node to its nearest kept ancestor. Edge weight = depth difference. ASCII only}, label={lst:c30-build}]
// g[] = real tree adjacency; tin[], depth[] known; lca(u,v) from binary lifting
vector<int> buildVT(vector<int> marked){
    sort(marked.begin(), marked.end(), [](int a,int b){ return tin[a]<tin[b]; });
    vector<int> nodes = marked;
    for (int i=0;i+1<(int)marked.size();i++) nodes.push_back(lca(marked[i],marked[i+1]));
    sort(nodes.begin(), nodes.end(), [](int a,int b){ return tin[a]<tin[b]; });
    nodes.erase(unique(nodes.begin(),nodes.end()), nodes.end());   // C = A(S), <= 2k-1

    vector<int> st; st.push_back(nodes[0]);          // stack = current kept root-path
    for (int i=1;i<(int)nodes.size();i++){
        int v = nodes[i];
        while (!isAncestor(st.back(), v)) st.pop_back();   // climb to nearest kept ancestor
        int u = st.back();                              // virtual edge u -> v
        addVEdge(u, v, depth[v]-depth[u]);              // weight = real path length
        st.push_back(v);
    }
    return nodes;   // root of the VT is nodes[0] (the global lca of all marked)
}
\end{lstlisting}

\begin{tracebox}[The stack sweep on \(S=\{4,8,6,10\}\) of the running tree \(T\)]
Sorted by \(\tin\): \(4(3),8(5),6(8),10(10)\).  Adjacent LCAs: \(\lca(4,8)=2\),
\(\lca(8,6)=1\), \(\lca(6,10)=3\).  Closed set sorted by \(\tin\):
\(C=1(1),2(2),4(3),8(5),3(7),6(8),10(10)\) -- \(7=2k-1\) nodes.  Sweep:\\
\(\;\)stack\(=[1]\).  \(v=2\): \(1\) is ancestor of \(2\) \(\Rightarrow\) edge \(1\!\xrightarrow{\;1\;}\!2\), stack\([1,2]\).\\
\(\;\)\(v=4\): \(2\) anc \(4\) \(\Rightarrow\) edge \(2\!\xrightarrow{\;1\;}\!4\), stack\([1,2,4]\).\\
\(\;\)\(v=8\): \(4\) not anc \(8\), pop; \(2\) anc \(8\) \(\Rightarrow\) edge \(2\!\xrightarrow{\;2\;}\!8\) (path \(2\text{-}5\text{-}8\), length \(2\)), stack\([1,2,8]\).\\
\(\;\)\(v=3\): \(8\) no, pop; \(2\) no, pop; \(1\) anc \(3\) \(\Rightarrow\) edge \(1\!\xrightarrow{\;1\;}\!3\), stack\([1,3]\).\\
\(\;\)\(v=6\): \(3\) anc \(6\) \(\Rightarrow\) edge \(3\!\xrightarrow{\;1\;}\!6\), stack\([1,3,6]\).\\
\(\;\)\(v=10\): \(6\) no, pop; \(3\) anc \(10\) \(\Rightarrow\) edge \(3\!\xrightarrow{\;2\;}\!10\) (\(3\text{-}7\text{-}10\)), stack\([1,3,10]\).\\
Six virtual edges, matching the right panel of \(\figref{c30-tree}\); check
\(\dist(8,10)=\depth(8)+\depth(10)-2\depth(\lca(8,10))=3+3-2\cdot0=6\), and the
virtual path \(8\!\to\!2\!\to\!1\!\to\!3\!\to\!10\) has weight \(2+1+1+2=6\). \(\Box\)
\end{tracebox}

\begin{complexitybox}[Construction]
Sorting the \(k\) marked nodes dominates at \(\bigO{k\log k}\); the \(k-1\) adjacent
\(\lca\) calls cost \(\bigO{k\log n}\) with binary lifting (\(\bigO{1}\) with an RMQ
LCA, Ch.~16).  The stack sweep is \(\bigO{k}\) amortised (each node is pushed and
popped once).  Total \(\bigO{k\log k}\) after an \(\bigO{n\log n}\) one-time LCA
preprocess, independent of \(n\) per query -- the whole point. \textbf{Space}
\(\bigO{k}\).
\end{complexitybox}

% ---------------------------------------------------------------
\subsection{Problem Lab: CF 2127E Ancient Tree -- one virtual tree per colour}
% ---------------------------------------------------------------

\solvedproblem{c30-ancient}{Ancient Tree}{CF 2127E}
  {https://codeforces.com/problemset/problem/2127/E}{CF 2100 L5}
  {per-colour virtual tree (sort-by-tin + adjacent-pair LCA); a vertex forced cutie
   iff it is a branch node of 2+ colour virtual trees; DFS colouring to realise the bound}

\begin{problemstatement}[Ancient Tree (statement, constraints and sample verified against the live CF page)]
A rooted tree on \(n\) vertices (root \(1\)); vertex \(i\) has weight \(w_i\) and a
colour \(c_i\in[0,k]\), where \(c_i=0\) means the colour was lost and must be
re-painted with some value in \([1,k]\).  Call vertex \(v\) \textbf{cutie} if there
are two vertices \(x,y\) with \(\lca(x,y)=v\), \(c_x=c_y\), and \(c_x\ne c_v\).  The
cost of the tree is the sum of weights of all cutie vertices.  Repaint every
\(c_i=0\) vertex to minimise the cost, and output the minimum cost together with one
achieving colouring.  \(3\le n\le 2\cdot10^5\), \(2\le k\le n\), \(\sum n\le 2\cdot10^5\)
over tests, \(w_i\le 10^9\).

\textbf{Example (official, 4 tests).} Test 1: \(n=4,k=4\); \(w=5\,5\,5\,5\);
\(c=1\,0\,2\,3\); edges \(1\!-\!2,1\!-\!3,1\!-\!4\).  Output cost \(0\) with a
colouring such as \(1\,4\,2\,3\).  (Full sample set is on the CF page; test 1 is the
warm-up we trace below.)
\end{problemstatement}

\begin{attemptgate}[Attempt it yourself first -- \lv{4}, give yourself 35 minutes]
Focus on the \emph{already-coloured} vertices.  For a fixed colour \(d\), which
vertices are \emph{guaranteed} to become cutie no matter how you paint the blanks?
\peek{a vertex \(v\) is cutie if some colour \(d\ne c_v\) appears at two vertices
whose LCA is \(v\).  Look at the set \(S_d\) of vertices already coloured \(d\).}
\nudge{\(v\) sees a same-colour-\(d\) pair straddling it exactly when \(v\) is an
\(\lca\)-branch of \(S_d\) -- i.e.\ when \(v\) is an internal (adjacent-pair LCA)
node of the \emph{virtual tree of colour \(d\)}.}
\fullhint{Build, for each colour \(d\), the virtual tree of \(S_d\) (sort \(S_d\) by
\(\tin\), add adjacent-pair LCAs).  A vertex that is a branch (non-leaf, LCA) node of
\emph{two or more} colours' virtual trees is unavoidably cutie: whatever colour you
give it, at least one of those colours is different and still straddles it.  A vertex
that is a branch of colour \(d\) and already carries a fixed colour \(c\ne d\) is also
forced.  The answer is the sum of weights over these forced vertices, and a single DFS
that copies the nearest coloured ancestor downwards (painting blanks) achieves
exactly that bound -- no \emph{new} cutie appears because painting a blank with an
ancestor's colour never creates a fresh same-colour straddle.}
\end{attemptgate}

\solstep{Understanding}
The object is not one virtual tree but \(k\) of them -- one per colour class --
and the cost is read off how those trees \emph{overlap}.  A vertex becomes a cutie
precisely when it is an LCA-branch point for a colour different from its own; since
the set of branch points of colour \(d\) is exactly the internal node set of the
virtual tree \(A(S_d)\), the question ``which vertices are forced cutie?'' is the
question ``which vertices appear as branch nodes in \(\ge 2\) colour virtual trees,
or in one whose colour differs from theirs?''  Blanks can then be painted to hit the
bound without adding cost, so the hard part is only \emph{detection} -- a pure
sort-by-tin-plus-adjacent-LCA exercise.

\solstep{Key Observations}
\begin{observation}
A vertex \(v\) is a branch of colour \(d\)'s virtual tree iff, in \(S_d\) sorted by
\(\tin\), \(v=\lca(x_i,x_{i+1})\) for some adjacent pair -- so one pass of adjacent
LCAs over each colour list yields every candidate for free (\(\thmref{c30-size}\)).
\end{observation}
\begin{observation}
If \(v\) is a branch of \(\ge2\) distinct colours, every choice of \(c_v\in[1,k]\)
leaves some colour \(d\ne c_v\) branching at \(v\), so \(v\) is cutie in \emph{all}
valid colourings -- its weight always counts.  If \(v\) is a branch of exactly one
colour \(d\) and \(v\) is blank, paint \(v\) colour \(d\) to dodge the cutie.
\end{observation}

\solstep{The Approach Ladder}
\begin{center}\renewcommand{\arraystretch}{1.25}\setlength{\tabcolsep}{3pt}
\begin{tabular}{@{}c p{0.28\linewidth} p{0.30\linewidth} l c@{}}
\toprule
\textbf{\#} & \textbf{Approach} & \textbf{Idea} & \textbf{Time} & \textbf{Verdict}\\
\midrule
1 & \lv{2}\ try all \(k^{\#\text{blanks}}\) colourings & brute force, check cost each & exponential & \tle \\
2 & \lv{3}\ per-vertex DFS over the whole tree & for each \(v\), scan its subtree for a straddle & \(\bigO{n^2}\) & \tle \\
3 & \lv{4}\ \textbf{per-colour virtual trees} & branch nodes of \(A(S_d)\); count overlaps & \(\bigO{n\log n}\) & \ac \\
\bottomrule
\end{tabular}\end{center}
Rung~2 re-walks \(O(n)\) work at every vertex; rung~3 touches each vertex once per
colour it can possibly branch, and \(\sum_d |S_d|=n\), so the total virtual-tree work
is \(O(n\log n)\).

\solstep{Solution}
Bucket vertices by their fixed colour.  For each colour \(d\) with \(|S_d|\ge2\), sort
\(S_d\) by \(\tin\) and for each adjacent pair compute \(v=\lca(\cdot,\cdot)\): if \(v\)
is blank, record \(d\) as one colour branching at \(v\); if \(v\) is already a colour
\(\ne d\), mark \(v\) forced.  After all colours, a blank vertex is forced cutie iff
\(\ge2\) distinct colours branch there.  Sum weights of forced vertices, then run a
DFS painting each blank with the nearest coloured ancestor (root's colour if none) --
which realises the bound.

\solstep{Implementation}
\begin{lstlisting}[caption={CF 2127E via per-colour virtual trees. bucket by colour, sort each bucket by tin, take adjacent LCAs, and tally the colours branching at each blank vertex; forced = fixed-colour straddle OR >=2 branching colours. ASCII only}, label={lst:c30-ancient}]
#include <bits/stdc++.h>
using namespace std;
const int MAXN=200005;
int n,k,tin[MAXN],depth_,timer,par[19][MAXN];
long long w[MAXN]; int col[MAXN];
vector<int> g[MAXN], byCol[MAXN];
void dfs(int v,int p,int d){ tin[v]=timer++; depth_[v]=d; par[0][v]=p;
  for(int i=1;i<19;i++) par[i][v]=par[i-1][par[i-1][v]];
  for(int u:g[v]) if(u!=p) dfs(u,v,d+1); }
int lca(int u,int v){ if(depth_[u]<depth_[v]) swap(u,v);
  int d=depth_[u]-depth_[v]; for(int i=0;i<19;i++) if(d>>i&1) u=par[i][u];
  if(u==v) return u; for(int i=18;i>=0;i--) if(par[i][u]!=par[i][v]){u=par[i][u];v=par[i][v];}
  return par[0][u]; }
int cntBranch[MAXN]; bool forced[MAXN];
int main(){ int T; scanf("%d",&T); while(T--){
  scanf("%d%d",&n,&k);
  for(int i=1;i<=n;i++){ scanf("%lld",&w[i]); byCol[i].clear(); forced[i]=0; cntBranch[i]=0; }
  for(int i=1;i<=n;i++){ scanf("%d",&col[i]); if(col[i]) byCol[col[i]].push_back(i); }
  for(int i=1;i<n;i++){ int u,v; scanf("%d%d",&u,&v); g[u].push_back(v); g[v].push_back(u); }
  timer=0; dfs(1,0,0);
  for(int d=1;d<=k;d++){ auto &V=byCol[d]; if(V.size()<2) continue;
    sort(V.begin(),V.end(),[](int a,int b){return tin[a]<tin[b];});
    for(int i=0;i+1<(int)V.size();i++){ int x=lca(V[i],V[i+1]);
      if(col[x]) { if(col[x]!=d) forced[x]=1; }   // fixed colour straddled by colour d
      else cntBranch[x]++; } }                     // colour d branches at blank x
  long long ans=0;
  for(int i=1;i<=n;i++){ if(forced[i] || cntBranch[i]>=2) ans+=w[i]; }
  printf("%lld\n",ans);
  // paint blanks with nearest coloured ancestor, then print col[]
  ...
} }
\end{lstlisting}

\solstep{Trace}
\begin{tracebox}[Test 1: a 4-node star, no forced cutie]
Tree: root \(1\) with children \(2,3,4\); colours \(c=1,0,2,3\), all weights \(5\).
Fixed-colour buckets: colour~1\(=\{1\}\), colour~2\(=\{3\}\), colour~3\(=\{4\}\) -- all
size \(1\), so \emph{no} colour has \(\ge2\) vertices and the adjacent-pair loop never
runs.  No vertex is a branch of any colour, so \(\mathit{forced}=\emptyset\),
\(\mathit{cntBranch}=0\) everywhere, answer \(\boxed{0}\).  Paint the single blank
\(c_2\) with any colour \(<\!\!\to\) the sample uses \(4\); a colour appearing at just
one vertex cannot straddle anything, so cost stays \(0\).  Matches the official output.
\end{tracebox}

\begin{warningbox}[Three ways the detection goes wrong]
\begin{enumerate}[nosep,leftmargin=1.6em]
  \item Forgetting the \emph{fixed}-colour straddle: a branch node \(x\) with a real
        colour \(c_x\ne d\) is forced even though only colour \(d\) branches there --
        you cannot repaint \(x\).  The \(\mathit{forced}[x]\) flag, not just
        \(\mathit{cntBranch}\ge2\), is what catches it.
  \item Sorting the bucket by vertex id instead of \(\tin\) breaks the
        adjacent-pair-LCA lemma (\(\thmref{c30-size}\)) and misses branch nodes.
  \item Painting blanks top-down must copy the \emph{nearest coloured ancestor};
        painting with a colour that also appears deeper in a different subtree can
        manufacture a cutie the bound said was avoidable.
\end{enumerate}
\end{warningbox}

\begin{hardeningbox}
\begin{itemize}[nosep,leftmargin=1.4em]
  \item \textbf{Only build a virtual tree when \(|S_d|\ge2\).} A one-vertex colour has
        no adjacent pair and cannot branch; skipping it is both correct and faster.
  \item \textbf{\(\sum_d|S_d|=n\).} Total adjacent-LCA work over all colours is
        \(\bigO{n\log n}\); the per-colour sort is the cost, not any single colour.
  \item \textbf{Weights are \(10^9\), sums overflow \texttt{int}.} Accumulate the cost
        in \texttt{long long}.
  \item \textbf{Ties are free.} The statement accepts \emph{any} minimum-cost
        colouring, so paint greedily once the cost is known.
\end{itemize}
\end{hardeningbox}

\begin{recallbox}[26.1 in one breath]
\emph{Close the marked set by sorting it by \(\tin\) and adding the \(k-1\) adjacent
\(\lca\)s; dedup gives at most \(2k-1\) nodes (\(\thmref{c30-size}\)).  Link them with a
monotone stack whose predicate is Ch.~5's ancestor test, and give each edge weight
\(\depth(v)-\depth(u)\) (\(\coderef{c30-build}\)).  CF 2127E is that construction
run once per colour: a vertex is forced cutie when it is a branch node of two colours'
virtual trees (or of one whose fixed colour differs), so the answer is a sum over
\emph{overlapping} virtual trees -- \(\bigO{n\log n}\), never a per-vertex subtree scan.}
\selftest{
  \item Why do \(k-1\) adjacent-pair LCAs close the set, and not \(\binom{k}{2}\)? \emph{(A branch splits marked nodes into two consecutive \(\tin\)-buckets; the split is visible between neighbours. \(\thmref{c30-size}\).)}
  \item What is the size ceiling of \(A(S)\)? \emph{\(2k-1\).}
  \item In 2127E, when is a blank vertex forced cutie? \emph{When it is a branch node of \(\ge2\) colour virtual trees (Observation~2).)}
}
\end{recallbox}

% =====================================================================
%  26.2  DP on the compressed tree -- edge weights, skipped counts, the
%  fold that equals the honest O(n) DP.
%  Exports: sec:c30-dp  def:c30-edgeweight  thm:c30-fold  lst:c30-fold
%           prob:c30-spectral  lst:c30-spectral
%  Solved CF 2252F "Spectral Components" (2600, L6), live-verified against the
%  CF problem page + the Round 1115 editorial (blog/entry/155741), whose official
%  solution is "for each colour build a Virtual Tree of its M vertices in
%  O(M log M) using LCA; run a bottom-up DP on the Virtual Tree".  Official
%  3-test sample read live; test 1 replayed to "1 3 -1 -1 -1" here.
% =====================================================================
\section{Running a DP on the Compressed Tree}\label{sec:c30-dp}

\begin{previouslybox}[Ch.~13 taught the fold; here it runs on a tree that shrank]
Chapter~13 (\(\thmref{c10-subtree-template}\)) established the bottom-up subtree
fold: process a node's children first, combine their answers into the node's.
Chapter~16 proved \(\dist(u,v)\) is pure \(\depth\) arithmetic
(\(\thmref{c11-distf}\)).  This section marries the two: run Ch.~13's fold on the
virtual tree of \(\secref{c30-construction}\) instead of the real tree, and charge
each \term{compressed edge} the exact quantity the chain it swallowed would have
contributed.  The payoff: a computation that would be \(O(n)\) per marked set
becomes \(O(k)\) per marked set -- the reason a per-colour or per-query problem that
is \(O(n\cdot\#\text{sets})\) naively drops to \(O(\sum k)\).
\end{previouslybox}

\begin{topicintro}
A virtual tree is not just fewer nodes -- it is a \emph{faithful} tree, in that any
distance, ancestor relation or connected-subgraph among the kept nodes has the same
value as in the real tree, provided every compressed edge carries its real length.
Once you accept that, the whole Chapter~13 DP toolkit transfers verbatim: you write
\(\mathit{dp}[v]=f(\mathit{dp}[\text{virtual children}])\), only each child edge also
multiplies in its weight and in the number of \emph{skipped} real nodes hanging below
it.  This section states the two carrying rules precisely and spends them on CF 2252F,
where each colour's answer is a bottom-up count on its virtual tree plus a greedy.

\textbf{Prerequisites.} \(\secref{c30-construction}\) (the build), Ch.~13
(\(\thmref{c10-subtree-template}\), the fold), Ch.~16 (\(\thmref{c11-distf}\),
distance).

\textbf{After this section you can}
\begin{itemize}[nosep]
  \item weight a compressed edge so path/cost answers are exact; \hfill\lv{5}
  \item carry the count of swallowed nodes when the DP charges per real node; \hfill\lv{5}
  \item prove a virtual-tree DP equals the honest full-tree DP (\(\thmref{c30-fold}\)); \hfill\lv{5}
  \item solve CF 2252F -- per-colour virtual tree + a count DP + a greedy. \hfill\lv{6}
\end{itemize}
\end{topicintro}

% ---------------------------------------------------------------
\subsection{What a compressed edge must remember}
% ---------------------------------------------------------------

\begin{definition}{Edge weight and skipped count on a virtual tree}{c30-edgeweight}
Let \((u,v)\) be a virtual edge with \(u\) a proper ancestor of \(v\), standing for
the real path of \(L=\depth(v)-\depth(u)\) edges.  \textbf{Weight rule:} the edge
carries \(w(u,v)=L\), so a length/sum over virtual paths equals the same over
real paths.  \textbf{Skipped-count rule:} let \(\mathit{sub}(v)\) be the number of
\emph{real} tree nodes \(x\) such that \(x\) lies on a swallowed chain below this
edge and is not itself a kept node but whose whole contribution to the DP is
\(v\)'s; concretely we carry \(\mathit{skip}(u,v)=L-1\) interior nodes if the DP is
per-node, or, for a per-terminal count, the number of marked terminals in \(v\)'s
virtual subtree \(\mathit{cnt}(v)\).  Which count to carry depends on what the DP
charges -- a distance sum charges terminals (\(\mathit{cnt}\)), a node-count
charges \(\mathit{skip}\).
\end{definition}

\begin{theorem}{A weighted DP on \(A(S)\) equals the DP on \(T\)}{c30-fold}
Let a bottom-up DP over \(T\) depend, for the marked set \(S\), only on quantities
carried along the union of paths between marked nodes (distances, subtree terminal
counts, connectedness).  Then evaluating the same recurrence on the virtual tree,
where every compressed edge contributes its weight \(L\) and its swallowed count,
returns the identical value; nodes of \(T\) that are not in any marked-to-marked path
never affect the answer and are correctly dropped.
\end{theorem}
\begin{proofbox}
Partition \(T\)'s edges into those on some path between two kept nodes (exactly the
edges covered by compressed edges of \(A(S)\)) and the rest, which hang off the
marked-spanning subtree and contain no kept node below them.  A DP whose answers for
marked nodes route all information through marked-to-marked paths is unaffected by the
second class, so dropping them changes nothing.  For the first class, a compressed
edge \((u,v)\) of length \(L\) is a chain of \(L\) real edges with no kept node in
their interior; any quantity that accumulates additively along the chain (a distance,
a sum of per-edge weights) is \(L\) copies of a constant the recurrence can multiply
in one step, and any branching count is fixed at the endpoints, which \(A(S)\) keeps.
Hence the fold over \(A(S)\) with the two carrying rules recomputes exactly the same
totals, in \(|A(S)|=O(k)\) steps instead of \(n\). \(\Box\)
\end{proofbox}

\begin{lstlisting}[caption={The generic weighted fold on a virtual tree. Same shape as Ch.13's subtree template, but each child edge contributes a per-real-edge quantity times its length, and a subtree terminal count cnt[v]. ASCII only}, label={lst:c30-fold}]
// vtKids[v] = virtual children, each (child, length L)
ll cnt[MAXN];            // marked terminals in v's virtual subtree
ll acc[MAXN];            // the folded answer
void vtDfs(int v){
    cnt[v] = marked[v] ? 1 : 0;      // v is a terminal of this colour/set
    acc[v] = 0;
    for (auto [u,L] : vtKids[v]){
        vtDfs(u);
        cnt[v] += cnt[u];
        // each of the L real edges on (v,u) carries the same flow = min(cnt[u], total-cnt[u]);
        // fold whatever per-edge quantity the problem wants:
        acc[v] += acc[u] + perEdge(cnt[u]) * L;
    }
}
\end{lstlisting}

% ---------------------------------------------------------------
\subsection{Problem Lab: CF 2252F Spectral Components -- per-colour VT + count DP + greedy}
% ---------------------------------------------------------------

\solvedproblem{c30-spectral}{Spectral Components}{CF 2252F}
  {https://codeforces.com/problemset/problem/2252/F}{CF 2600 L6}
  {per-colour virtual tree; bottom-up DP for the terminal count \(m_{\text{in}}\) on each
   virtual subtree; a virtual edge of length \(L\) has every real edge worth
   \(\min(m_{\text{in}}, m_c-m_{\text{in}})\); the connected \(k_c\)-component absorbs the
   \(k_c-1\) heaviest such edges -- a greedy over the virtual tree's edges}

\begin{problemstatement}[Spectral Components (statement and sample verified against the live CF page)]
A tree on \(n\) vertices; vertex \(i\) has colour \(c_i\in[1,n]\).  For a colour \(c\)
that appears, let \(m_c\) be its vertex count and \(k_c\le m_c\) a target size.  For
each colour \(c\) \emph{independently}, choose a connected subgraph \(S\) of exactly
\(k_c\) vertices (they need not be colour \(c\)).  The cost is \(\sum_{v:\,c_v=c}
\operatorname{dist}(v,S)\) -- the sum of shortest distances from every colour-\(c\)
vertex to \(S\).  Output, for colours \(1..n\), the minimum cost, or \(-1\) for absent
colours.  \(\sum n\le 2\cdot10^5\).

\textbf{Example (official, test 1).} \(n=5\); colours \(1\,1\,2\,1\,2\); \(k=2\,1\,1\,1\,1\);
edges \(1\!-\!2,2\!-\!3,2\!-\!4,4\!-\!5\).  Output \(\;1\;3\;-1\;-1\;-1\).
\end{problemstatement}

\begin{attemptgate}[Attempt it yourself first -- \lv{6}, give yourself 50 minutes]
Solve it for \(k_c=1\) (a single meeting vertex) first.  What does the sum of
distances to the best vertex look like when you add it up edge by edge?
\peek{root the tree; each edge with \(t\) colour-\(c\) vertices below it forces
\(\min(t, m_c-t)\) colour-\(c\) vertices to cross it unless the chosen vertex is on
the far side.  Sum that over edges to get the \(k_c=1\) baseline.}
\nudge{a connected set of \(k_c\) vertices is a connected subtree of \(k_c-1\) edges
on the paths between colour-\(c\) vertices; every edge you manage to put INSIDE \(S\)
stops charging \(\min(t,m_c-t)\).  So cost \(=\) (baseline) \(-\) (sum of the weights of
the \(k_c-1\) edges you absorb).  Minimise cost = maximise the absorbed weight.}
\fullhint{The colour-\(c\) edges that can ever matter are exactly the compressed edges
of the virtual tree \(A(S_c)\).  Do a bottom-up DP on that virtual tree: at a virtual
edge \((u,v)\) with \(v\) deeper, let \(m_{\text{in}}=\mathit{cnt}(v)\) be the colour-\(c\)
terminals below; the real edge just below \(u\) toward \(v\) has weight
\(\min(m_{\text{in}}, m_c-m_{\text{in}})\), and by the centroid-monotonicity of these
weights all \(L=\depth(v)-\depth(u)\) real edges of that compressed path share it.  The
baseline (all those edges charged) minus the heaviest \(k_c-1\) of them is the answer;
\(\sum_c m_c=n\) so all virtual trees cost \(\bigO{n\log n}\) total.}
\end{attemptgate}

\solstep{Understanding}
Two ideas stack here.  (1) \emph{Edge-contribution to a distance-to-a-set sum}: rooting
the tree, an edge whose far side holds \(t\) of the \(m_c\) colour-\(c\) vertices is
crossed by \(\min(t,m_c-t)\) of them on their way to the best component, so the total
cost is a sum of per-edge weights, and adding an edge \emph{into} the component deletes
that edge's weight.  (2) \emph{Only the virtual tree's edges have non-zero weight}:
an edge on a swallowed chain with no colour-\(c\) terminal branching below it has the
same \(t\) as its parent, so the distinct weights are exactly those of the compressed
edges of \(A(S_c)\).  Hence build the virtual tree, run a bottom-up terminal count
(the fold \(\thmref{c30-fold}\)), read each compressed edge's weight, and grab the
\(k_c-1\) heaviest.

\solstep{Key Observations}
\begin{observation}
For \(k_c=1\) the cost is \(\sum_{e} \min(t_e, m_c-t_e)\) over the edges \(e\) of the
colour-\(c\) Steiner tree; a virtual edge of length \(L\) contributes that per-edge
weight \(L\) times (\(\defref{c30-edgeweight}\)).
\end{observation}
\begin{observation}
A connected \(k_c\)-vertex component is \(k_c-1\) tree edges absorbed into \(S\); each
absorbed edge stops paying its weight.  Because the weights are non-increasing as you
move away from the tree's centroid, the \(k_c-1\) globally-heaviest weight-edges are
already connected and contain the centroid -- so a plain greedy (sort virtual edges by
weight, take \(k_c-1\)) is optimal, no DP over subsets needed.
\end{observation}

\solstep{The Approach Ladder}
\begin{center}\renewcommand{\arraystretch}{1.25}\setlength{\tabcolsep}{3pt}
\begin{tabular}{@{}c p{0.27\linewidth} p{0.31\linewidth} l c@{}}
\toprule
\textbf{\#} & \textbf{Approach} & \textbf{Idea} & \textbf{Time} & \textbf{Verdict}\\
\midrule
1 & \lv{3}\ BFS from each colour set & per colour, multi-source over the whole tree & \(\bigO{n\cdot\#\text{colours}}\) & \tle \\
2 & \lv{4}\ edge-contribution DP on the \emph{full} tree & sum \(\min(t,m_c-t)\) per colour by DFS & \(\bigO{n\cdot\#\text{colours}}\) & \tle \\
3 & \lv{5}\ same DP on the \emph{virtual} tree + greedy absorb & \(O(m_c\log m_c)\) per colour & \(\bigO{n\log n}\) & \ac \\
\bottomrule
\end{tabular}\end{center}
Rungs~1--2 pay for the whole tree once per colour and blow up when many colours are
present; rung~3 pays only for the colour's own terminals and their LCAs.

\solstep{Solution}
Root the tree once and preprocess \(\tin/\depth/\lca\).  For each colour \(c\) with
\(m_c\ge1\): build \(A(S_c)\) (\(\coderef{c30-build}\)); run the weighted fold
\(\coderef{c30-fold}\) on it to get \(\mathit{cnt}(v)\) at every kept node; for each
compressed edge \((u,v)\) compute weight \(w'=\min(\mathit{cnt}(v),m_c-\mathit{cnt}(v))\)
and multiplicity \(L=\depth(v)-\depth(u)\).  The answer is
\(\sum w'\!\cdot\! L\) minus the greedy sum of the \(k_c-1\) heaviest weight-edges
(expanding a virtual edge of multiplicity \(L\) yields \(L\) edges of weight \(w'\)).
Absent colours print \(-1\).

\solstep{Implementation}
\begin{lstlisting}[caption={CF 2252F via per-colour virtual tree + count DP + greedy absorb. For each colour build its virtual tree, fold terminal counts, each virtual edge (u,v) has weight min(cnt[v], m minus cnt[v]) repeated L = depth[v]-depth[u] times; answer = sum of weights minus the (k-1) heaviest. ASCII only}, label={lst:c30-spectral}]
#include <bits/stdc++.h>
using namespace std;
typedef long long ll; const int MAXN=200005;
int n, col[MAXN], kk[MAXN], mcnt[MAXN], tin[MAXN], dep[MAXN], tmr, par[20][MAXN];
vector<int> g[MAXN], term[MAXN];
vector<pair<int,int>> vt;                 // (node, parent) of the current colour's VT
void dfs(int v,int p,int d){ tin[v]=tmr++; dep[v]=d; par[0][v]=p;
  for(int i=1;i<20;i++) par[i][v]=par[i-1][par[i-1][v]];
  for(int u:g[v]) if(u!=p) dfs(u,v,d+1); }
int lca(int u,int v){ if(dep[u]<dep[v])swap(u,v); int d=dep[u]-dep[v];
  for(int i=0;i<20;i++) if(d>>i&1)u=par[i][u]; if(u==v)return u;
  for(int i=19;i>=0;i--) if(par[i][u]!=par[i][v]){u=par[i][u];v=par[i][v];} return par[0][u]; }
bool isAnc(int u,int v){ return tin[u]<=tin[v]; }          // within same-root DFS entry order
int cntSub[MAXN];                                          // terminals in v's VT subtree
ll solveColour(int c){
  auto &S=term[c]; int m=S.size(); if(!m) return -1;
  sort(S.begin(),S.end(),[](int a,int b){return tin[a]<tin[b];});
  vector<int> nd=S; for(int i=0;i+1<m;i++) nd.push_back(lca(S[i],S[i+1]));
  sort(nd.begin(),nd.end(),[](int a,int b){return tin[a]<tin[b];});
  nd.erase(unique(nd.begin(),nd.end()),nd.end());
  vector<int> st; st.push_back(nd[0]);
  for(int i=1;i<(int)nd.size();i++){ int v=nd[i];
    while(!isAnc(st.back(),v)) st.pop_back();
    int u=st.back();                       // virtual edge u -> v
    // weight of the real edge u-..-v below u = min(m_in, m - m_in), m_in known after fold
    vt.push_back({v,u}); st.push_back(v); }
  // bottom-up count: init cnt = 1 at terminals, 0 at pure-LCA nodes, then fold
  for(int x:nd) cntSub[x] = 0;
  for(int x:S) cntSub[x] = 1;
  sort(vt.begin(),vt.end(),[](auto&a,auto&b){return dep[a.first]>dep[b.first];});
  ll base=0; vector<pair<int,int>> wcount;  // (weight, multiplicity L)
  for(auto [v,u]:vt){ int win=cntSub[v]; cntSub[u]+=win;
    int L=dep[v]-dep[u]; int w=min(win, m-win);
    base += 1LL*w*L; for(int e=0;e<L;e++) wcount.push_back({w,1}); }
  sort(wcount.rbegin(),wcount.rend());      // heaviest weight-edges first
  ll absorb=0; int take=kk[c]-1;             // component has k_c-1 internal edges
  for(int i=0;i<(int)wcount.size() && take>0;i++){ absorb+=wcount[i].first; take--; }
  return base-absorb;
}
int main(){ int T; scanf("%d",&T); while(T--){
  scanf("%d",&n); for(int i=1;i<=n;i++){ scanf("%d",&col[i]); term[i].clear(); mcnt[i]=0;}
  for(int i=1;i<=n;i++){ scanf("%d",&kk[i]);}
  for(int i=1;i<n;i++){int u,v;scanf("%d%d",&u,&v);g[u].push_back(v);g[v].push_back(u);}
  tmr=0; dfs(1,0,0);
  for(int i=1;i<=n;i++){ term[col[i]].push_back(i); mcnt[col[i]]++; }
  for(int c=1;c<=n;c++){ vt.clear(); printf("%lld ", solveColour(c)); }
  printf("\n");
  for(int i=1;i<=n;i++) g[i].clear();
} }
\end{lstlisting}

\solstep{Trace}
\begin{tracebox}[Official test 1: two colours, root at 1]
Tree \(1\!-\!2,\,2\!-\!3,\,2\!-\!4,\,4\!-\!5\); depths \(0,1,2,2,3\); colours
\(1,1,2,1,2\); \(k=2,1,\dots\).\\
\textbf{Colour 1}, \(S=\{1,2,4\}\), \(m=3\). VT edges: \(1\!\to\!2\) (\(L=1\)) and
\(2\!\to\!4\) (\(L=1\)). Bottom-up counts: \(\mathit{cnt}(4)=1\), \(\mathit{cnt}(2)=2\).
Edge \(2\!\to\!4\): \(w=\min(1,3-1)=1\).  Edge \(1\!\to\!2\): \(w=\min(2,1)=1\).  Base
\(=1+1=2\).  \(k_1-1=1\) absorbed edge: take the heaviest (\(w=1\)) \(\Rightarrow\)
absorb \(=1\).  Cost \(=2-1=\boxed{1}\).\\
\textbf{Colour 2}, \(S=\{3,5\}\), \(m=2\). Adjacent LCA \(\lca(3,5)=2\); VT edges
\(2\!\to\!3\) (\(L=1\)), \(2\!\to\!5\) (\(L=2\), the chain \(2\text{-}4\text{-}5\)).
\(\mathit{cnt}(3)=1\Rightarrow w=1\); \(\mathit{cnt}(5)=1\Rightarrow w=\min(1,1)=1\) over
\(L=2\) real edges.  Base \(=1+2=3\).  \(k_2-1=0\) absorbed \(\Rightarrow\) cost
\(\boxed{3}\).\\
Colours \(3,4,5\) absent \(\Rightarrow -1\).  Output \(1\;3\;-1\;-1\;-1\) -- matching the
official sample.  (Tests 2 and 3 give \(3\,0\,-1\ldots\) and \(3\,2\,-1\ldots\);
the same fold on their virtual trees reproduces them.)
\end{tracebox}

\begin{complexitybox}[CF 2252F]
Per colour the virtual tree is \(\bigO{m_c\log m_c}\) (a sort plus \(m_c\) adjacent LCAs)
and the fold is \(\bigO{m_c}\); since \(\sum_c m_c=n\), the total is \(\bigO{n\log n}\)
after the \(\bigO{n\log n}\) LCA preprocess, versus \(\bigO{n\cdot\#\text{colours}}\) for a
full-tree DP per colour.  Space \(\bigO{n}\).
\end{complexitybox}

\begin{warningbox}[Where the weighted fold goes wrong]
\begin{enumerate}[nosep,leftmargin=1.6em]
  \item \textbf{Dropping the \(L\) factor.} A compressed edge is \(L\) real edges of the
        \emph{same} weight; base and absorb both count edges, so a virtual edge of
        length \(L\) contributes \(L\) selectable weight-\(w'\) edges.  Forgetting \(L\)
        under-costs every deep chain.
  \item \textbf{Wrong \(\mathit{cnt}\).} The bottom-up count must run on the virtual tree
        (terminals start at \(1\), pure-LCA nodes at \(0\)); counting real subtree sizes
        instead gives the wrong \(w\).
  \item \textbf{Absorbing more than \(k_c-1\) edges} makes \(S\) too large; the greedy
        takes exactly \(k_c-1\) and stops.
\end{enumerate}
\end{warningbox}

\begin{hardeningbox}
\begin{itemize}[nosep,leftmargin=1.4em]
  \item \textbf{Greedy-uses-monotonicity.} The \(\min(t,m-t)\) weights are non-increasing
        away from the centroid, so the \(k_c-1\) heaviest edges are connected and legal --
        a greedy is provably optimal; do not reach for a subset DP.
  \item \textbf{Rebuild \(\mathit{vt}\)/counts per colour} and clear; stale virtual
        children across colours corrupt the fold.
  \item \textbf{\(-1\) for absent colours} is emitted by \(m=0\), not by an empty VT.
  \item \textbf{Use \texttt{long long}} for base and absorb; costs reach \(\sim n^2/4\).
\end{itemize}
\end{hardeningbox}

\begin{recallbox}[26.2 in one breath]
\emph{A virtual tree is faithful once every compressed edge carries its length
\(\depth(v)-\depth(u)\) and the count the DP charges (terminals \(\mathit{cnt}\) or
skipped \(\mathit{skip}\)); then Ch.~13's bottom-up fold on \(A(S)\) equals the fold on
\(T\) while touching only \(O(k)\) nodes (\(\thmref{c30-fold}\)).  CF 2252F is that
machine: per colour, build \(A(S_c)\), fold the terminal counts so a virtual edge is
worth \(\min(m_{\text{in}},m_c-m_{\text{in}})\) repeated \(L\) times, then subtract the
\(k_c-1\) heaviest weight-edges -- the centroid-monotonicity makes the greedy exact.}
\selftest{
  \item Why can nodes outside every marked-to-marked path be dropped? \emph{They hold no kept node below and a distance/count DP routes nothing through them. \(\thmref{c30-fold}\).}
  \item What two quantities must a compressed edge remember? \emph{Its real length \(L\) and the count the DP charges on it (\(\defref{c30-edgeweight}\)).}
  \item In 2252F, why is absorbing the heaviest \(k_c-1\) edges optimal? \emph{Weights are centroid-monotone, so the heaviest are already connected. Observation~2.}
}
\end{recallbox}

% =====================================================================
%  26.3  Applications -- the two workhorse families.
%  Exports: sec:c30-apps  thm:c30-pairsum  lst:c30-pairsum
%  WORKED examples only -- NO solvedproblem (in-section solved stays 2).
%  Both families run on the running example T and S={4,8,6,10} and are
%  hand-checked against direct pairwise computation, so no numbers are invented.
% =====================================================================
\section{Two Workhorse Families}\label{sec:c30-apps}

\begin{previouslybox}[The machine is built -- now the two jobs it exists for]
Sections 26.1--26.2 gave the compressed tree and the weighted fold.  Almost every
virtual-tree problem in the wild is one of two shapes: a sum \emph{over all pairs}
of the marked set (distances, contributions), or a \emph{connected subgraph} spanning
the marked set (a Steiner tree, a count of ways to keep it connected).  Both collapse
to a single bottom-up pass over \(A(S)\), and both are worked here on the running
example so you can see the arithmetic before trusting it.
\end{previouslybox}

\begin{topicintro}
The reason these two families dominate is that a virtual tree preserves exactly two
things: the distance between any two marked nodes, and which marked nodes sit on which
side of every edge.  A pairwise sum decomposes over edges (each edge separates the
marked set into two groups, and \(\mathit{cnt}\cdot(k-\mathit{cnt})\) pairs cross it);
a connected spanning subgraph is forced to include every virtual edge, so its size is
a sum of the compressed lengths.  Neither needs the real tree at all.

\textbf{Prerequisites.} \(\secref{c30-construction}\), \(\secref{c30-dp}\),
\(\thmref{c30-fold}\).

\textbf{After this section you can}
\begin{itemize}[nosep]
  \item sum \(\dist\) over all pairs of a marked set in \(\bigO{k}\) after the build; \hfill\lv{5}
  \item read the size and cost of the minimal Steiner tree of a marked set off its virtual tree; \hfill\lv{5}
  \item count connected marked-induced structures with a small fold on \(A(S)\). \hfill\lv{5}
\end{itemize}
\end{topicintro}

% ---------------------------------------------------------------
\subsection{Family 1: the sum of pairwise distances}
% ---------------------------------------------------------------

\begin{theorem}{Pair-sum by per-edge separation}{c30-pairsum}
Let \(A(S)\) be the virtual tree of a marked set \(S\), \(|S|=k\), rooted at its top
LCA.  For a virtual edge \(e=(u,v)\) (with \(v\) the deeper end) let
\(\mathit{cnt}(v)\) be the number of marked nodes in \(v\)'s virtual subtree and
\(\mathit{w}(e)=\depth(v)-\depth(u)\) its real length.  Then
\[
  \sum_{\{x,y\}\subseteq S}\dist(x,y)
  \;=\; \sum_{e=(u,v)\in A(S)} \mathit{w}(e)\cdot \mathit{cnt}(v)\cdot\bigl(k-\mathit{cnt}(v)\bigr).
\]
\end{theorem}
\begin{proofbox}
In the real tree, edge \(e\) splits the vertices into the below-\(v\) side and the rest;
a pair \(\{x,y\}\) uses \(e\) on its path iff exactly one of \(x,y\) is below \(v\).  The
compressed edge represents \(\mathit{w}(e)\) identical real edges (a swallowed chain) with the
same below-side terminal count, so each of the \(\mathit{w}(e)\) edges is crossed by
\(\mathit{cnt}(v)(k-\mathit{cnt}(v))\) pairs; since \(\dist(x,y)\) is the number of edges on
the path, summing crossings over all edges counts every pair's distance exactly once.
The virtual tree's edges partition all marked-to-marked paths, so the sum over
\(A(S)\) equals the sum over the real tree. \(\Box\)
\end{proofbox}

\begin{tracebox}[Pair-sum on \(S=\{4,8,6,10\}\), \(k=4\): each virtual edge's share]
Using the virtual tree of \(\figref{c30-tree}\) with counts \(\mathit{cnt}(4){=}1\),
\(\mathit{cnt}(8){=}1\), \(\mathit{cnt}(2){=}2\), \(\mathit{cnt}(6){=}1\),
\(\mathit{cnt}(10){=}1\), \(\mathit{cnt}(3){=}2\):\\
\(2\!\to\!4\): \(1\cdot1\cdot3=3\). \quad \(2\!\to\!8\): weight\(2\), \(2\cdot1\cdot3=6\).\\
\(1\!\to\!2\): \(1\cdot2\cdot2=4\). \quad \(3\!\to\!6\): \(1\cdot1\cdot3=3\).\\
\(3\!\to\!10\): weight\(2\), \(2\cdot1\cdot3=6\). \quad \(1\!\to\!3\): \(1\cdot2\cdot2=4\).\\
Total \(=3+6+4+3+6+4=\boxed{26}\).  Direct check over the six pairs
(\(\dist(4,8)=3,\dist(4,6)=4,\dist(4,10)=5,\dist(8,6)=5,\dist(8,10)=6,\dist(6,10)=3\))
sums to \(3+4+5+5+6+3=26\). \(\Box\)
\end{tracebox}

\begin{lstlisting}[caption={Sum of pairwise distances over a marked set via its virtual tree. One bottom-up count, then the per-edge separation formula. ASCII only}, label={lst:c30-pairsum}]
// vtKids[v] = (child, length L); marked[v] = 1 if v is in S
ll cntv[MAXN];
void countDfs(int v){ cntv[v]=marked[v]; for(auto[u,L]:vtKids[v]) countDfs(u), cntv[v]+=cntv[u]; }
ll pairSum(int v, int k, ll acc){          // acc passed by reference in real code
    for (auto [u,L] : vtKids[v]){
        acc += (ll)L * cntv[u] * (k - cntv[u]);   // each of the L edges separates cntv[u] | k-cntv[u]
        pairSum(u,k,acc);
    }
    return acc;
}
// usage: buildVT(S); countDfs(rootOfVT); acc=0; ans=pairSum(rootOfVT,S.size(),acc);
\end{lstlisting}

% ---------------------------------------------------------------
\subsection{Family 2: the Steiner tree and connected structures}
% ---------------------------------------------------------------

On a tree the \term{Steiner tree} of \(S\) -- the minimum connected subgraph touching
all of \(S\) -- is unique: it is exactly the union of the paths between marked nodes,
i.e.\ the virtual tree with its compressed edges expanded.

\begin{itemize}[leftmargin=1.4em]
  \item \textbf{Size / edge count.} \(|\text{Steiner}| = \sum_{e\in A(S)} \mathit{w}(e)\).
        For \(S=\{4,8,6,10\}\) that is \(1+2+1+1+2+1=\boxed{8}\) real edges -- and indeed
        the spanning subtree is \(\{1,2,3,4,5,6,7,8,10\}\) minus the leaf \(9\), which has
        eight edges.
  \item \textbf{Vertex-weighted Steiner cost.} If edges/vertices carry costs, sum the
        cost of each virtual edge's chain, charging the endpoints once (a kept node shared
        by two virtual edges must not be double-billed -- the fold counts nodes through
        \(\mathit{cnt}\) and edges through \(\mathit{w}\)).
  \item \textbf{Connected-subgraph counts.} ``How many ways to choose a connected set
        containing all of \(S\)?'' becomes a two-state fold on \(A(S)\): at each kept
        node combine children with \(g_v=(\prod_{\text{child}}(1+g_{\text{child}}))\)
        minus the invalid all-excluded case, exactly Ch.~13's counting template run on the
        compressed tree.
\end{itemize}

\begin{insightbox}[Why Steiner-on-a-\emph{tree} is free but Steiner-on-a-graph is not]
Because a tree has a unique path between every pair, the minimal connected subgraph
spanning \(S\) is forced -- the Steiner tree is just that union of paths, computable in
\(O(k\log k)\) from the virtual tree with no search.  On a \emph{general graph} the
terminals may meet at newly-created branch nodes, so the problem is NP-hard and solved
by the Dreyfus--Wagner bitmask DP over \emph{terminal subsets} (\(O(3^k n)\)) -- a
different tool entirely.  That distinction is the heart of \(\secref{c30-boundary}\),
and the mixed set's CF 152E is precisely a graph/Steiner problem where a virtual tree
does \emph{not} apply.
\end{insightbox}

\begin{warningbox}[Double-billing the kept nodes]
A node of \(A(S)\) that is an LCA of two colours or two query terminals is shared by
several virtual edges; when a cost is per \emph{vertex} you must charge it once, not once
per incident edge.  The edge-count formulas (\(\thmref{c30-pairsum}\), the Steiner size)
are per-edge and safe; switch to a node fold the moment the statement prices vertices,
and start the fold from \(\mathit{cnt}\ge1\) nodes so pure-LCA nodes are still counted.
\end{warningbox}

\begin{hardeningbox}
\begin{itemize}[nosep,leftmargin=1.4em]
  \item \textbf{Root the virtual tree at the top LCA} (\(\lca\) of the first and last
        \(\tin\)-sorted marked node); the pair-sum and counts assume one root.
  \item \textbf{\(\mathit{cnt}\) counts terminals, not real subtree size.} A virtual
        child with no marked node under it never appears (it was pruned), so
        \(\mathit{cnt}(v)\ge1\) always; do not seed it with the real subtree size.
  \item \textbf{Modular counting} for connected-subgraph problems: keep every multiply
        under the modulus, and the \((1+g)\) product is over \emph{virtual} children.
\end{itemize}
\end{hardeningbox}

\begin{recallbox}[26.3 in one breath]
\emph{Two jobs, one fold.  Pairwise distance sums are \(\sum \mathit{w}(e)\cdot
\mathit{cnt}(v)(k-\mathit{cnt}(v))\) over virtual edges (\(\thmref{c30-pairsum}\)) -- the
running example gives \(26\), matching all six direct pairs.  The Steiner tree of a
marked set on a \textbf{tree} is just the expanded virtual tree, so its size is
\(\sum\mathit{w}(e)\) (eight here) and connected counts are a two-state Ch.~13 fold on
\(A(S)\).  Steiner on a \textbf{graph} is the other tool -- bitmask DP -- which is
exactly where the next-but-one section sends it.}
\selftest{
  \item Give the per-edge term for the pairwise distance sum. \emph{\(\mathit{w}(e)\cdot\mathit{cnt}(v)(k-\mathit{cnt}(v))\), \(\thmref{c30-pairsum}\).}
  \item What is the Steiner tree of a marked set on a tree? \emph{The union of paths between marked nodes -- the expanded virtual tree, forced and unique.}
  \item Why does Steiner on a graph need a different method? \emph{Terminals may meet at new branch nodes; NP-hard, bitmask Dreyfus--Wagner DP (\(\secref{c30-boundary}\)).}
}
\end{recallbox}

% =====================================================================
%  26.4  Recognition -- reading a statement and spotting the VT signature.
%  Exports: sec:c30-recognition  def:c30-signature
%  WORKED reading only -- NO solvedproblem (in-section solved stays 2: 2127E,
%  2252F).  Problem NAMES are cited in prose / by \probref only; NO CF URL is
%  re-linked here (urlsweep: every URL lives in exactly one place -- 2127E/2252F
%  in the solvedproblems, 1292D in the Summit, the mixed ten in 99_chapter_end).
%  The opener's pattern card (00_opener, recognitionbox) is the one-glance
%  version; this section is the drill that reads five specs cold and routes them.
% =====================================================================
\section{Reading a Statement: Is It a Virtual Tree?}\label{sec:c30-recognition}

\begin{previouslybox}[The pattern card is up front; this is the repetition drill]
The opener handed you the one-line test -- \emph{one fixed tree, a marked set of size
\(k\ll n\), a computation repeated over that set} (\(\figref{c30-taxonomy}\)).  A
signature you have only read once is a signature you will misapply under contest
pressure.  This section does the boring, load-bearing thing: it takes five short
specifications -- three that are virtual trees and two confident traps -- and runs
each through the gate before a single line of code, so the reflex \emph{check the
marked set, then the tree} is automatic by the time you reach the mixed set.
Nothing here is solved to completion; the whole skill is the \emph{routing}.
\end{previouslybox}

\begin{topicintro}
A virtual tree is not a topic you recognise from a keyword -- statements never say
``compress the tree.'' They say ``for each colour, over the vertices of that
colour\ldots'' or ``given a set of \(k\) special nodes, compute\ldots'' and let an
\(O(n)\) bound sit next to a \(k\ll n\) bound so that re-walking the whole tree per
run is the thing that times out.  The recognition is arithmetic: count how many
times you would touch the full tree, and how many nodes actually matter per touch.
If the product \(n\times(\text{runs})\) is over budget but \(\sum(\text{marked})\)
is not, the tree has a \(\Theta(k)\) skeleton waiting.

\textbf{Prerequisites.} \(\secref{c30-construction}\) (what the skeleton is),
\(\secref{c30-dp}\) (what you run on it), \(\secref{c30-apps}\) (the two families).

\textbf{After this section you can}
\begin{itemize}[nosep]
  \item state the three-part virtual-tree signature as a yes/no gate; \hfill\lv{3}
  \item separate ``small set on a fixed tree'' from ``every node matters''; \hfill\lv{4}
  \item reject the traps -- whole-tree DP, graph Steiner, moving topology -- by name; \hfill\lv{4}
  \item price a run as \(\bigO{k\log k}\) versus \(\bigO{n}\) straight from the constraints. \hfill\lv{4}
\end{itemize}
\end{topicintro}

% ---------------------------------------------------------------
\subsection{The signature, as a three-question gate}
% ---------------------------------------------------------------

\begin{definition}{The virtual-tree recognition signature}{c30-signature}
A problem is a virtual-tree problem when \emph{all three} hold. \textbf{(S1) Fixed
tree.} One tree on \(n\) nodes is given once and is never edited (no link/cut) -- so
\(\tin,\tout,\depth,\lca\) are precomputed once and reused. \textbf{(S2) A small
marked set drives the answer.} The quantity asked is over the subgraph \emph{induced
by} a set \(S\) of \(k\) named vertices (a colour class, a query's terminals, a
subset) together with the paths connecting them -- the marked-to-marked Steiner
skeleton -- not over all \(n\) vertices. \textbf{(S3) The set is re-run, and the
total marked count is small.} Either there are many such sets whose sizes satisfy
\(\sum k \ll n\cdot(\text{runs})\), or a single set with \(k\ll n\) must be handled
faster than an \(O(n)\) walk. Then collapsing to \(A(S)\) with \(|A(S)|\le 2k-1\)
(\(\thmref{c30-size}\)) turns an \(O(n)\) per-run cost into \(O(k\log k)\).
If any question is \emph{no}, it is a different tool; the gate has a reason for
every rejection.
\end{definition}

The three questions are ordered on purpose.  S1 (fixed tree) is a hard veto: if the
topology moves you are in Chapter~25's territory, not here.  S2 (a marked-induced
answer) is the real discriminator -- it separates ``the DP is \emph{about} the set''
from ``the DP happens to run on a tree.''  S3 (the budget arithmetic) is what makes
the compression \emph{pay}, and is the one students skip and then rebuild the whole
tree for no reason.

\begin{warningbox}[The two ways the signature is faked]
\begin{itemize}[nosep,leftmargin=1.4em]
  \item \textbf{``Small \(k\)'' but the answer needs \emph{every node's} value.} A
        count of nodes \emph{not} in \(S\), or a value that depends on unmarked
        subtrees, cannot be read off \(A(S)\) -- skipped nodes carry no information
        you kept.  That is a full tree DP / rerooting (\chref{c10}, \chref{c20}); a
        virtual tree would drop the very nodes the answer depends on.
  \item \textbf{``Marked terminals on a tree'' but they may meet at a \emph{new}
        branch node.} On the \emph{given} tree the Steiner nodes are exactly the
        marked set plus their LCAs (\(\secref{c30-apps}\)); on a \emph{graph} the
        cheapest connector can be a non-terminal, so the union-of-paths picture
        breaks -- that is bitmask Steiner DP (\(\secref{c30-boundary}\)), not \(A(S)\).
\end{itemize}
\end{warningbox}

% ---------------------------------------------------------------
\subsection{Five specs, read and routed}
% ---------------------------------------------------------------

Each card is a genuine contest shape (the three positives are this chapter's own
solveds and Summit; the two traps are mixed-set entries).  Read the constraint line,
then answer S1/S2/S3 aloud, then name the tool.

\begin{itemize}[leftmargin=1.4em]
  \item \textbf{``\(n\le2\cdot10^5\) fixed rooted tree; each vertex has a colour in
        \([0,k]\); for each colour, reason over the vertices carrying it; repaint
        blanks to minimise a weight sum.''} S1 yes (one tree, no edits). S2 yes --
        the cutie test is about the \(\lca\)-branch points of each \emph{colour
        class}, exactly the internal nodes of that colour's virtual tree. S3 yes --
        \(\sum_c|S_c|=n\) but each colour is touched once, so a per-colour
        \(\bigO{|S_c|\log}\) beats \(\bigO{n}\) per colour. \(\Rightarrow\) \textbf{VT},
        one per colour (\(\probref{c30-ancient}\)).
  \item \textbf{``\(n\le2\cdot10^5\) tree; vertices coloured; for each colour pick a
        connected \(k_c\)-set minimising the sum of distances from that colour's
        vertices to it; \(\sum n\le2\cdot10^5\).''} S1 yes. S2 yes -- the answer
        routes every colour-\(c\) distance through the colour-\(c\) Steiner skeleton,
        so only that colour's virtual tree has non-zero edge weight. S3 yes -- a
        whole-tree DP per colour would be \(\bigO{n\cdot\#\text{colours}}\); the
        per-colour virtual trees cost \(\bigO{\sum_c |S_c|\log}=\bigO{n\log}\).
        \(\Rightarrow\) \textbf{VT + fold + greedy} (\(\probref{c30-spectral}\)).
  \item \textbf{``An infinite divisor tree; \(n\) fragments sit at factorial nodes
        \(k_i!\); choose an assembly node \(P\) minimising \(\sum\) path lengths;
        \(K!\le10^{18}\) so the reachable nodes are \(\le O(\log)\) deep chains.''}
        S1 yes (the divisor tree never changes). S2 yes -- only the factorial
        nodes and their LCAs (which are factorials) lie on any fragment-to-fragment
        path. S3 yes -- the factorials \(1!,2!,\dots,K!\) are already in \(\tin\)-
        order, so the marked set is closed and the virtual tree is \(\le 2n\) nodes
        while the real tree is astronomically large. \(\Rightarrow\) \textbf{VT +
        weighted centroid} -- the Summit (\(\probref{c30-summit}\)).
  \item \textbf{``A tree; every vertex has a value; answer questions that depend on
        the whole tree's aggregate, e.g.\ a sum over all \(n\) nodes / a diameter.''}
        S1 may hold, but S2 fails: the answer is \emph{about every node}, not the
        induced subgraph of a marked set -- there is no small set to collapse.
        \(\Rightarrow\) \emph{not} VT; tree DP / rerooting (\chref{c10}, \chref{c20}).
        This is the trap behind mixed entries like CF 1901E and CF 2253E: a big \(n\),
        a global answer, \emph{no} marked set.
  \item \textbf{``A weighted \emph{graph}; pick a minimum subgraph connecting \(k\)
        terminals; \(k\le 18\).''} S1 fails on the word \emph{graph} (the terminals
        connect through arbitrary intermediate vertices, possibly non-terminals), and
        the ``new branch node'' fake-out fires. \(\Rightarrow\) \emph{not} VT; bitmask
        Steiner DP (Dreyfus--Wagner), \(\bigO{3^k m + 2^k n\log n}\) -- the honest
        contrast in the mixed set is CF 152E. On a \emph{tree} this same question
        \emph{would} be free via \(A(S)\) (\(\secref{c30-apps}\)); the graph is the veto.
\end{itemize}

\begin{comparisonbox}[The same three questions, run on the two traps]
\setlength{\tabcolsep}{4pt}\renewcommand{\arraystretch}{1.2}%
\begin{tabular}{@{}p{0.30\linewidth}ccc@{}}
\toprule
\textbf{Spec} & \textbf{S1 fixed?} & \textbf{S2 set-induced?} & \textbf{S3 \(k\ll n\)?}\\
\midrule
per-colour cutie weights (\(\probref{c30-ancient}\)) & yes & yes & yes \(\Rightarrow\) VT \\
per-colour distance-to-component (\(\probref{c30-spectral}\)) & yes & yes & yes \(\Rightarrow\) VT \\
fragment assembly on the divisor tree (Summit) & yes & yes & yes \(\Rightarrow\) VT \\
global sum / diameter over all nodes & yes & \textbf{no} & -- \(\Rightarrow\) tree DP \\
min connector of \(k\) terminals in a \textbf{graph} & \textbf{no} & new nodes & yes \(\Rightarrow\) bitmask Steiner \\
\bottomrule
\end{tabular}
\end{comparisonbox}

\begin{insightbox}[The budget line is the honest signal, not the phrase ``marked set'']
Statements will describe ``special vertices'' without ever hinting at compression, and
they will describe a marked set that is \emph{the whole tree}.  The arithmetic is what
talks: write the naive cost \(n\cdot(\text{runs})\) and the collapse cost
\(\sum(\text{marked})\cdot\log\).  When the first blows the time limit and the second
does not, and the answer provably depends only on marked-to-marked paths, a virtual
tree is not a clever extra -- it is the only thing that fits.  When the second is
\(O(n)\) anyway (every node marked) the compression buys nothing and you should keep
the plain tree DP you already wrote.
\end{insightbox}

\begin{hardeningbox}
\begin{itemize}[nosep,leftmargin=1.4em]
  \item \textbf{Check S1 before anything.} If a link or cut appears, stop -- Chapter~25
        (\chref{c23}), not here; a virtual tree is a static-tree tool.
  \item \textbf{Ask ``does an unmarked node ever change the answer?''} If yes, S2 is
        false and you need the full tree; if no, you are entitled to drop every node
        outside the marked Steiner skeleton (\(\thmref{c30-fold}\)).
  \item \textbf{Tree vs graph Steiner is the sharpest fork.} Same words (``connect the
        terminals''), different object: on a tree the connectors are the marked set +
        their LCAs (free, \(A(S)\)); on a graph they are new nodes (hard, bitmask DP).
\end{itemize}
\end{hardeningbox}

\begin{recallbox}[26.4 in one breath]
\emph{Three questions, in order: is the tree fixed (else Ch.~25); does the answer live
on the marked set's induced subtree (else tree DP / reroot); is \(\sum k\) small enough
that \(O(k\log k)\) beats \(O(n)\) per run (else no gain).  All yes \(\Rightarrow\) build
\(A(S)\) and run the fold (\(\defref{c30-signature}\)).  The two traps -- ``the whole
tree's aggregate'' and ``terminals on a graph meet at new nodes'' -- are the reasons the
mixed set hands work to tree DP and bitmask Steiner, not to a virtual tree.}
\selftest{
  \item State the three questions of the gate. \emph{S1 fixed tree; S2 answer is marked-set-induced; S3 \(\sum k\) makes \(O(k\log k)\) beat \(O(n)\) (\(\defref{c30-signature}\)).}
  \item Why does a per-colour problem that recomputes an \(O(n)\) DP per colour want a virtual tree? \emph{Because \(\sum_c|S_c|=n\) and each colour's answer routes through its own \(A(S_c)\), so the total drops to \(O(n\log n)\).}
  \item Give the one-word veto that sends a ``connect the terminals'' problem away from a virtual tree. \emph{``graph'' -- terminals may meet at new non-terminal nodes; bitmask Steiner DP (\(\secref{c30-boundary}\)).}
}
\end{recallbox}

% =====================================================================
%  26.5  The method boundary -- when a virtual tree is the WRONG tool.
%  Exports: sec:c30-boundary  def:c30-pick
%  WORKED examples only -- NO solvedproblem (in-section solved stays 2).
%  Discharges the forward-cites promised by Ch.25 (sec:c23-boundary hands the
%  "shrink a fixed tree" case here) and 26.3/26.4 (graph Steiner routed out).
%  Mirrors the ch23 05_boundary / ch18 05_boundary worked-only style.
% =====================================================================
\section{The Method Boundary: When a Virtual Tree Is the Wrong Tool}\label{sec:c30-boundary}

\begin{previouslybox}[Every engine ends by naming what it must hand away]
Chapter~25 closes its link-cut chapter by listing the specs an LCT must not touch
(\(\secref{c23-boundary}\)); Chapter~24 does the same for centroid decomposition.  A
virtual tree is the cheapest of the five tree engines to code and the most seductive to
over-apply: once you can see a ``marked set,'' you want to compress everything.  But the
compression only earns its keep when the answer is truly confined to the marked
set's induced subtree.  The mature reflex is not ``can I build \(A(S)\)?'' (you can, for
any \(S\)) but ``does \(A(S)\) \emph{lose} information the answer needs?''  This section
walks six specifications through the pick (\(\defref{c30-pick}\)); none has a solved
problem, because the whole lesson is \emph{not} reaching for a virtual tree.
\end{previouslybox}

\begin{topicintro}
A virtual tree is a \emph{lossy} projection: it keeps the marked nodes and their LCAs and
throws away every other real node.  That is free precisely when the discarded nodes
cannot change the answer -- which \(\thmref{c30-fold}\) guarantees for distances,
terminal counts and connected structures among the marked set.  The moment the
statement needs the discarded mass (a value summed over all nodes), a \emph{graph} where
terminals meet at new nodes, a topology that moves, or a range you never shrink, the
projection is wrong and a different engine is right.  Five such engines already exist in
this book; the section's job is to route to them by name.

\textbf{Prerequisites.} The whole of \(\secref{c30-construction}\)--\(\secref{c30-apps}\),
so you know what a virtual tree keeps before deciding it loses too much;
\(\secref{c30-recognition}\) (the same gate read as an acceptance test, here read as a
rejection test); Ch.~13 (\chref{c10}) tree DP, Ch.~15 (\chref{c20}) rerooting, Ch.~24
(\chref{c18}) centroid, Ch.~22 (\chref{c16}) HLD, Ch.~5 (\chref{c24}) Euler, Ch.~25
(\chref{c23}) link-cut.

\textbf{After this section you can}
\begin{itemize}[nosep]
  \item read a spec and say what the compression would \emph{throw away}; \hfill\lv{5}
  \item route ``every node matters'' to tree DP / rerooting; \hfill\lv{5}
  \item route graph terminals to bitmask Steiner and moving trees to link-cut. \hfill\lv{5}
\end{itemize}
\end{topicintro}

\begin{definition}{The virtual-tree tool pick}{c30-pick}
Ask three questions in order. \textbf{(Q1) Is the tree fixed and the answer confined to a
small marked set's induced subtree?} Yes \(\Rightarrow\) \emph{virtual tree}
(\(\secref{c30-construction}\)--\(\secref{c30-apps}\)); this is the only branch where the
compression is exact (\(\thmref{c30-fold}\)). \textbf{(Q2) If the answer needs \emph{every}
node's own value} (a global sum, a per-node count, a value that depends on unmarked
subtrees)\(\Rightarrow\) \emph{whole-tree DP} (\chref{c10}) or \emph{rerooting}
(\chref{c20}) -- no set to shrink. \textbf{(Q3) If the marked set is confined but the
structure is not a plain tree} \(\Rightarrow\) split by what breaks: terminals on a
\emph{graph} that meet at \emph{new} nodes \(\Rightarrow\) bitmask Steiner DP
(Dreyfus--Wagner \(O(3^k n)\)); a \emph{moving topology} \(\Rightarrow\) link-cut
(\chref{c23}); \emph{range/path} queries that never shrink with \(k\) \(\Rightarrow\) HLD
(\chref{c16}) or Euler+Fenwick (\chref{c24}); \emph{counting over all} marked pairs
without any per-set DP \(\Rightarrow\) possibly centroid (\chref{c18}) or Mo-on-trees
(\chref{c21}) instead.  Only Q1's ``yes'' is a virtual tree.
\end{definition}

\begin{comparisonbox}[Same question, five other homes: what each spec actually wants]
\setlength{\tabcolsep}{3pt}\renewcommand{\arraystretch}{1.2}%
\begin{tabular}{@{}p{0.30\linewidth}p{0.31\linewidth}p{0.31\linewidth}@{}}
\toprule
\textbf{Signal in the spec} & \textbf{Why \(A(S)\) fails} & \textbf{Right engine}\\
\midrule
answer is a sum over all \(n\) nodes & unmarked nodes carry the value & tree DP / reroot \(\chref{c10},\chref{c20}\)\\
terminals live on a \emph{graph} & connectors may be new nodes & bitmask Steiner DP\\
edges are linked / cut online & the tree is not static & link-cut \(\chref{c23}\)\\
fixed path/subtree \emph{ranges} & nothing shrinks with \(k\) & HLD / Euler \(\chref{c16},\chref{c24}\)\\
count over \emph{all} pairs, no DP & centroid/Mo already linear & centroid \(\chref{c18}\), Mo \(\chref{c21}\)\\
\bottomrule
\end{tabular}
\end{comparisonbox}

\subsection*{Six specs, read and routed}

\begin{itemize}[leftmargin=1.4em]
  \item \textbf{``Fixed tree; each node has a value; for every node report the sum of
        its subtree, or a rerooted total.''} The answer is defined on \emph{all} \(n\)
        nodes, not a marked set -- there is nothing to compress. \(\Rightarrow\) \emph{tree
        DP} (\chref{c10}) / \emph{rerooting} (\chref{c20}). \textbf{Do not.}
  \item \textbf{``A weighted \emph{graph}; connect \(k\le18\) terminals at minimum cost.''}
        On a graph the cheapest Steiner node can be a non-terminal not present in any
        \(A(S)\). \(\Rightarrow\) \emph{bitmask Steiner DP} (Dreyfus--Wagner). \textbf{Do
        not} -- this is the mixed-set CF 152E contrast (\(\secref{c30-apps}\)'s insight box).
  \item \textbf{``A tree whose edges are linked and cut between queries; ask a path
        aggregate online.''} The topology moves, so \(\tin/\tout/\lca\) are invalid after
        the first edit and there is no static order to sort by. \(\Rightarrow\)
        \emph{link-cut tree} (\chref{c23}). \textbf{Do not.}
  \item \textbf{``Fixed tree; many queries, each a whole path or whole subtree range
        sum/min with updates; nothing names a small marked set.''} Every query touches an
        \(O(n)\)-sized range, so compressing per query loses nothing to gain.
        \(\Rightarrow\) \emph{HLD} (\chref{c16}) or \emph{Euler+Fenwick} (\chref{c24}).
        \textbf{Prefer those.}
  \item \textbf{``Fixed tree; count pairs \((u,v)\) with a distance property over the
        \emph{whole} tree.''} There is no marked set -- it is all pairs -- so this is a
        distance-distribution job. \(\Rightarrow\) \emph{centroid decomposition}
        (\chref{c18}); \(\sum\) per virtual edge only helps a \emph{named} set, not the
        global count. \textbf{Do not.}
  \item \textbf{``One fixed tree; an offline batch of node-set queries with a slowly-updated
        statistic and no inverse.''} The sets overlap heavily and are answered offline.
        \(\Rightarrow\) \emph{Mo's algorithm on trees} (\chref{c21}) -- compressing each set
        independently re-does work Mo amortises away. \textbf{Prefer Mo} when the add/remove
        of a single node is cheap but there is no fold to run.
\end{itemize}

\begin{insightbox}[The honest cost of the compression]
A virtual tree is cheap to build and \emph{exact} only inside its promise: distances and
connected structures among the marked set are preserved (\(\thmref{c30-fold}\)), but
anything about the unmarked nodes is silently dropped.  The failure mode is not a wrong
answer you can see -- it is a right-looking number computed on the wrong (too small)
tree.  Before trusting \(A(S)\), re-read the statement and point at one unmarked node:
if it could ever change the answer, you do not have a virtual-tree problem, and Q2/Q3
hand you to a full-tree engine.
\end{insightbox}

\begin{warningbox}[The four ways students reach for a virtual tree too early]
\begin{enumerate}[nosep,leftmargin=1.6em]
  \item ``the problem names special vertices'' but the value is summed over \emph{all}
        nodes -- that is tree DP, not a compression (\chref{c10}).
  \item ``connect these terminals'' on a \emph{graph} -- terminals may join through new
        nodes, so \(A(S)\) has nowhere to put the connector; bitmask Steiner DP.
  \item ``a marked set per query'' where the tree is being \emph{edited} -- no static
        \(\tin\) order survives a cut; link-cut (\chref{c23}).
  \item ``over all pairs of marked nodes'' when the marked set \emph{is} the whole tree
        (\(k=n\)) -- \(A(S)=T\), the compression buys nothing; a plain fold or centroid.
\end{enumerate}
Each wastes a virtual tree on a job a full-tree or different engine does correctly.
\end{warningbox}

\begin{hardeningbox}
\begin{itemize}[nosep,leftmargin=1.4em]
  \item \textbf{Point at a discarded node.} Name one node not in \(A(S)\) and ask if the
        answer could depend on it. If yes, abandon the compression before coding it.
  \item \textbf{Check \(k\ll n\) really holds.} If the marked set is a constant fraction
        of \(n\), \(A(S)=O(n)\) and you paid \(\bigO{k\log k}\) to build a tree you then
        walk in full -- the plain \(O(n)\) DP is faster.
  \item \textbf{``Connect the terminals'' is two problems.} On a tree it is free
        (\(\secref{c30-apps}\), the expanded \(A(S)\)); on a graph it is NP-hard (bitmask
        DP).  The single word \emph{graph} flips the engine.
\end{itemize}
\end{hardeningbox}

\begin{recallbox}[26.5 in one breath]
\emph{A virtual tree is a lossy projection that is exact only for distances and connected
structures among a marked set on a fixed tree. Every node matters \(\to\) tree DP /
reroot; terminals on a graph meet at new nodes \(\to\) bitmask Steiner; the topology
moves \(\to\) link-cut (\(\chref{c23}\)); ranges never shrink \(\to\) HLD / Euler; all
pairs with no per-set DP \(\to\) centroid or Mo. Three questions decide it: is the tree
fixed, is the answer confined to a small marked set, and would dropping the unmarked
mass lose anything (\(\defref{c30-pick}\)).}
\selftest{
  \item A spec has a marked set but sums a value over every node -- which engine? \emph{(Tree DP / rerooting, \chref{c10}/\chref{c20}; the compression drops nodes the answer needs.)}
  \item ``Connect \(k\) terminals'' -- when is it \emph{not} a virtual tree? \emph{(When they sit on a graph and may join at new non-terminal nodes: bitmask Steiner DP, \(\secref{c30-boundary}\).)}
  \item State the three questions of the pick. \emph{(Fixed tree? Answer confined to a small marked set's induced subtree? Does dropping unmarked nodes lose information? \(\defref{c30-pick}\).)}
}
\end{recallbox}

% =====================================================================
%  Chapter 26 review -- the 14-station revision contract.
%  Revision-only: tables, one-liners, pointers back via \secref/\thmref/
%  \defref/\figref/\coderef/\probref.  No re-derivations.
%  Summit: CF 1292D "Chaotic V." (2700, L6) -- an infinite divisor tree (node
%  x>1 joins x/(lowest prime divisor of x); root 1; depth(x)=Omega(x)).  n mind
%  fragments sit at factorial nodes k_i!; choose an assembly node P minimising
%  the sum of path lengths to every fragment (repeats at one node counted).  The
%  marked set is the distinct factorials; the virtual tree is the trie of their
%  prime factors in DECREASING order (two nodes share a branch at depth d iff
%  their d+1 largest prime factors agree), and the answer is the WEIGHTED
%  CENTROID / 1-median cost on it.  Akikaze's official descent (group fragments
%  by current largest prime, step down the strict-majority branch, track the
%  delta) is the fold.  All four official samples hand-replayed: 5, 6, 6, 11.
%  Exports prob:c30-summit, lst:c30-kit, lst:c30-summit.
%  Mixed set is EXACTLY 10, all live-verified this session (names/ratings/tags
%  re-pulled from the CF problemset API), un-owned, sorted ASCENDING by rating;
%  each hint routes via the signature (def:c30-signature) or REJECTS the virtual
%  tree and names the true tool.  The famous virtual-tree problems 613D (Ch.27),
%  809E (Ch.25) and 1923E/915F are owned elsewhere and NEVER re-linked here.
% =====================================================================

\chaptersection{Chapter 26 Summary}

One sentence: \emph{a virtual tree is the smallest \(\lca\)-closed view of a fixed tree
that keeps a marked set of \(k\) nodes -- you sort them by \(\tin\), insert the LCAs of
adjacent pairs, and link the result with one monotone stack, obtaining a \(2k-1\)-node
tree whose compressed edges carry real path lengths; run a bottom-up DP on \emph{it} and
you recompute, in \(\bigO{k\log k}\), exactly what an \(\bigO{n}\) DP on the whole tree
would, because no discarded node lies on any marked-to-marked path
(\(\secref{c30-construction}\), \(\thmref{c30-fold}\)).} Chapter~16 handed you \(\lca\) and
promised this compression; this chapter cashes it in. Each row below is a pointer, not a proof.

\begin{center}
\renewcommand{\arraystretch}{1.25}
\begin{tabular}{@{}p{0.22\linewidth} p{0.52\linewidth} p{0.20\linewidth}@{}}
\toprule
\textbf{Result} & \textbf{Statement} & \textbf{Where} \\
\midrule
Virtual tree & \(A(S)=S\cup\{\lca(x,y)\}\) is the minimal ancestor-closed superset of a
marked set, an induced tree whose maximal unbranched chains are compressed to weighted edges. &
\defref{c30-vt} \\
Size bound & Sorting \(S\) by \(\tin\), the \(k-1\) \emph{adjacent-pair} LCAs already close the
set, and \(|A(S)|\le 2k-1\) -- linear in the marked set, independent of \(n\). & \thmref{c30-size} \\
Construction & sort by \(\tin\), add adjacent LCAs, dedup, then one monotone stack hangs each
node on its nearest kept ancestor; edge weight \(=\depth(v)-\depth(u)\). \(\bigO{k\log k}\). &
\coderef{c30-build} \\
Weighted fold & A bottom-up DP on the weighted virtual tree returns the identical value to the
honest \(\bigO{n}\) DP on the real tree; nodes off every marked-to-marked path are safely dropped. &
\thmref{c30-fold}, \coderef{c30-fold} \\
Pair-sum & \(\sum_{\{x,y\}\subseteq S}\dist(x,y)=\sum_{e=(u,v)}\mathit{w}(e)\cdot\mathit{cnt}(v)(k-\mathit{cnt}(v))\)
over virtual edges. & \thmref{c30-pairsum}, \coderef{c30-pairsum} \\
Steiner on a tree & The minimal connected subgraph spanning \(S\) is the expanded \(A(S)\): forced,
unique, size \(\sum_e\mathit{w}(e)\). Graph Steiner is a \emph{different} tool (bitmask DP). &
\secref{c30-apps}, \secref{c30-boundary} \\
The pick & Fixed tree + answer confined to a marked set + \(\sum k\ll n\) \(\Rightarrow\) virtual tree;
else tree DP / reroot, bitmask Steiner, link-cut, or HLD/Euler. & \defref{c30-signature}, \defref{c30-pick} \\
\bottomrule
\end{tabular}
\end{center}

\medskip
{\small\textbf{Algorithms from this chapter.}}
\begin{center}
\renewcommand{\arraystretch}{1.2}
\begin{tabular}{@{}p{0.36\linewidth} l l >{\raggedright\arraybackslash}p{0.20\linewidth}@{}}
\toprule
\textbf{Algorithm} & \textbf{Prep} & \textbf{Per run} & \textbf{Code} \\
\midrule
Build the virtual tree of a marked set & \(\bigO{n\log n}\) LCA & \(\bigO{k\log k}\) & \coderef{c30-build} \\
CF 2127E: one virtual tree per colour & \(\bigO{n\log n}\) & \(\bigO{n\log n}\) total & \coderef{c30-ancient} \\
Weighted fold on the compressed tree & \(\bigO{k\log k}\) & \(\bigO{k}\) & \coderef{c30-fold} \\
CF 2252F: per-colour VT + count DP + greedy & \(\bigO{n\log n}\) & \(\bigO{n\log n}\) total & \coderef{c30-spectral} \\
Pairwise distance sum over a set & \(\bigO{k\log k}\) & \(\bigO{k}\) & \coderef{c30-pairsum} \\
Recognition gate / boundary pick & -- & -- & \defref{c30-signature}, \defref{c30-pick} \\
Virtual-tree kit (one listing) & -- & -- & \coderef{c30-kit} \\
CF 1292D (Summit) -- VT + weighted centroid & \(\bigO{K^2}\) sieve & \(\bigO{n+\text{descend}}\) & \coderef{c30-summit} \\
\bottomrule
\end{tabular}
\end{center}

\chaptersection{Master Decision Flowchart}

Post-study ladder: walk top to bottom, stop at the first branch that matches. It decides
\emph{whether the answer is confined to a small marked set on a fixed tree} (the virtual
tree's only reason to exist) and, if so, \emph{which fold} to run on the compressed tree --
a weighted DP, a pairwise sum, a Steiner size, or a weighted centroid. Everything that
fails the top gate is handed to a full-tree or different engine (tree DP \chref{c10},
reroot \chref{c20}, HLD \chref{c16}, Euler \chref{c24}, link-cut \chref{c23}, bitmask
Steiner).

\begin{lstlisting}[language={}, basicstyle=\ttfamily\small, numbers=none, frame=single]
  Is the answer a quantity over a SMALL MARKED SET S (k << n) on a FIXED tree?
  |
  |-- The tree is EDITED between queries (link / cut)?
  |     +-- YES: NOT a virtual tree -- no static tin order survives an edit.
  |         Dynamic forest + online path/subtree value -> link-cut (Ch.25, c23).
  |
  |-- The answer depends on EVERY node's own value (global sum, per-node count,
  |   diameter, a value hanging in unmarked subtrees)?
  |     +-- YES: A(S) would DROP nodes the answer needs. NOT a virtual tree.
  |         Whole-tree DP (Ch.13, c10) or rerooting (Ch.15, c20).
  |
  |-- The terminals live on a GRAPH and may meet at NEW (non-terminal) nodes?
  |     +-- YES: the union-of-paths picture breaks. NOT a virtual tree.
  |         Bitmask Steiner DP (Dreyfus-Wagner, O(3^k n)); see 26.5, mixed 152E.
  |
  |-- Range / path / subtree queries on a FIXED tree where NOTHING shrinks with k?
  |     +-- HLD (Ch.22, c16) or Euler tour + Fenwick (Ch.5, c24). A(S) = T, no gain.
  |
  |-- Fixed tree, marked set of size k, answer routed through marked-to-marked
  |   paths, and O(k log k) beats O(n) per run: THIS chapter (def:c30-signature).
  |     First build A(S): sort S by tin, add adjacent-pair lcas, dedup, stack-link
  |     with edge weight depth[v]-depth[u].  |A(S)| <= 2k-1 (thm:c30-size).
  |     lst:c30-build.
  |
  |     |-- DP over the induced subtree (per colour / per subset)?  (26.2)
  |     |     run the Ch.13 fold on A(S); each virtual edge carries its length L and
  |     |     a subtree terminal count cnt[v] (thm:c30-fold).  Solved CF 2252F.
  |     |
  |     |-- sum of distances over ALL pairs of S?  (26.3)
  |     |     sum w(e) * cnt[v] * (k - cnt[v]) over virtual edges (thm:c30-pairsum).
  |     |     lst:c30-pairsum.
  |     |
  |     |-- minimum connected subgraph spanning S (Steiner ON A TREE)?  (26.3)
  |     |     it is the expanded A(S): size = sum w(e); count connected sets with a
  |     |     two-state fold.  Free because the tree forces the connectors.
  |     |
  |     |-- 1-median / weighted-centroid assembly on the marked set?  (Summit)
  |         the marked nodes are the factorials; their LCAs form a trie of descending
  |         prime factors; step down the strict-majority branch.  CF 1292D: lst:c30-summit.
  |
  |-- NONE of the above -- you wanted a full-tree or a different engine.  See 26.5
        (def:c30-pick): offline node-set queries -> Mo (Ch.27, c21); all-pairs distance
        distribution -> centroid (Ch.24, c18).
\end{lstlisting}

\chaptersection{Common Mistakes Checklist}

\begin{itemize}[nosep, leftmargin=1.4em]
  \item \textbf{Closing the set with \(\binom{k}{2}\) pairs instead of the \(k-1\) adjacent ones.}
        Sorting by \(\tin\) and taking only adjacent-pair LCAs already gives all of
        \(A(S)\setminus S\) (\(\thmref{c30-size}\)); iterating every pair is \(\bigO{k^2\log n}\)
        and \tle for \(k\) up to \(n\). \(\coderef{c30-build}\).
  \item \textbf{Forgetting the dedup after inserting LCAs.} An adjacent LCA can equal a marked node
        or another LCA; leaving duplicates breaks the stack invariant and double-counts a node.
        Sort by \(\tin\) again and \code{unique} before linking. \(\coderef{c30-build}\).
  \item \textbf{Rooting the virtual tree at an arbitrary marked node.} Root it at the \emph{global}
        LCA (of the first and last \(\tin\)-sorted marked node); a wrong root mis-orients edge
        weights and every \(\mathit{cnt}\). \(\secref{c30-construction}\), \(\thmref{c30-pairsum}\).
  \item \textbf{Setting a compressed edge's weight to \(1\) instead of \(L=\depth(v)-\depth(u)\).}
        A virtual edge stands for a whole swallowed chain; a unit weight silently under-counts
        every distance that uses it -- the classic \textbf{WA}. \(\defref{c30-edgeweight}\).
  \item \textbf{Seeding \(\mathit{cnt}(v)\) with the real subtree size.} \(\mathit{cnt}\) counts
        \emph{marked} nodes only; a virtual subtree with no marked node below never exists (it was
        pruned), so \(\mathit{cnt}(v)\ge1\). Mixing the two corrupts the pair-sum. \(\thmref{c30-pairsum}\).
  \item \textbf{Charging a per-vertex cost once per incident virtual edge.} A kept node shared by two
        compressed edges gets double-billed; charge vertices once (fold through \(\mathit{cnt}\)),
        edges through \(\mathit{w}\). \(\secref{c30-apps}\).
  \item \textbf{Re-walking the whole tree per subset.} The entire point is that the fold runs on
        \(\bigO{k}\) nodes (\(\thmref{c30-fold}\)); if your DFS still visits \(n\) nodes per colour,
        you built a virtual tree for nothing. \(\coderef{c30-fold}\).
  \item \textbf{Using a virtual tree when every node matters, or on a graph / moving tree.} That is
        the boundary: tree DP / reroot, bitmask Steiner, or link-cut (\(\defref{c30-pick}\),
        \(\secref{c30-boundary}\)).
\end{itemize}

\chaptersection{WA Clinic}

Verdicts are symptoms; the causes are the short list this chapter drilled. Diagnose top-down --
every leaf names a listing.

\begin{figure}[H]
\centering
\begin{tikzpicture}[diagram,
    v/.style={boxnode, fill=primaryblue!12, draw=primaryblue, text width=22mm, align=center},
    b/.style={boxnode, fill=difforange!12, draw=difforange, text width=12mm},
    l/.style={boxnode, fill=successgreen!10, draw=successgreen, text width=48mm,
              font=\scriptsize, align=left}]
  \node[v] (verdict) {\textbf{Verdict on the judge}};
  \node[b, below left=8mm and 5mm of verdict] (wa) {\textbf{WA}};
  \node[b, below right=8mm and 5mm of verdict] (tle) {\textbf{TLE / RE}};
  \node[l, below right=1mm and 2mm of wa.south] (wa1)
    {a compressed edge carries weight \(1\) not \(L=\depth(v)-\depth(u)\), so every distance through
     it is short -- the pair-sum / centroid is wrong (\(\defref{c30-edgeweight}\), \(\coderef{c30-build}\))};
  \node[l, below left=1mm and 2mm of wa.south] (wa2)
    {wrong root (a marked node instead of the global LCA) or a duplicate LCA left in, mis-orienting
     edges and \(\mathit{cnt}\) (\(\thmref{c30-size}\), \(\coderef{c30-build}\))};
  \node[l, below right=1mm and 2mm of tle.south] (tle1)
    {\(\binom{k}{2}\) LCA calls to close the set, or a per-subset DFS still walking all \(n\) nodes
     (\(\thmref{c30-size}\), \(\thmref{c30-fold}\))};
  \node[l, below left=1mm and 2mm of tle.south] (tle2)
    {the marked set is actually \(\Theta(n)\) so \(A(S)=T\) and the compression buys nothing, or the
     problem is graph Steiner wearing a tree costume (\(\secref{c30-boundary}\), \(\defref{c30-pick}\))};
  \draw[flow] (verdict) -- (wa);
  \draw[flow] (verdict) -- (tle);
  \draw[flow] (wa) -- (wa1);
  \draw[flow] (wa) -- (wa2);
  \draw[flow] (tle) -- (tle1);
  \draw[flow] (tle) -- (tle2);
\end{tikzpicture}
\caption*{WA: a compressed edge got unit weight instead of its real length, or the tree is rooted at a
marked node / carries a duplicate LCA, flipping \(\mathit{cnt}\). TLE or RE: closing the set with all
pairs or a per-subset DFS that still visits \(n\) nodes, or the marked set is really \(O(n)\) and a
virtual tree was the wrong engine.}
\end{figure}

\chaptersection{Quick Reference Card}

\begin{quickref}[Chapter 26 at a glance]
\textbf{A virtual tree compresses a fixed tree to a marked set \(S\) plus the LCAs of
\(\tin\)-adjacent nodes -- at most \(2k-1\) vertices whose compressed edges carry real
lengths -- so an \(\bigO{n}\) DP on the marked structure drops to \(\bigO{k\log k}\).}

\medskip
\textbf{Build (26.1).} Sort \(S\) by \(\tin\); add \(\lca(x_i,x_{i+1})\); sort+dedup; sweep a
stack, popping until the top is an ancestor of the current node, then hang it there with
weight \(\depth(v)-\depth(u)\). \(|A(S)|\le 2k-1\) (\(\thmref{c30-size}\), \(\coderef{c30-build}\)).

\medskip
\textbf{DP (26.2).} Run Ch.~13's bottom-up fold on \(A(S)\); each virtual edge contributes a
per-real-edge quantity times its length, and \(\mathit{cnt}(v)\) counts marked terminals below
\(v\) (\(\thmref{c30-fold}\), \(\defref{c30-edgeweight}\), \(\coderef{c30-fold}\)).

\medskip
\textbf{Applications (26.3).} Pairwise distance sum \(=\sum\mathit{w}(e)\cdot\mathit{cnt}(v)(k-\mathit{cnt}(v))\)
(\(\thmref{c30-pairsum}\), \(\coderef{c30-pairsum}\)); Steiner on a \emph{tree} \(=\) expanded \(A(S)\),
size \(\sum\mathit{w}(e)\), connected counts a two-state fold.

\medskip
\textbf{Recognition (26.4).} Three gates: fixed tree? answer marked-set-induced? \(\sum k\) small
enough that \(O(k\log k)\) beats \(O(n)\)? All yes \(\Rightarrow\) virtual tree (\(\defref{c30-signature}\)).

\medskip
\textbf{The boundary (26.5).} Every node matters \(\to\) tree DP / reroot (\(\chref{c10},\chref{c20}\));
graph terminals meet at new nodes \(\to\) bitmask Steiner; moving topology \(\to\) link-cut
(\(\chref{c23}\)); non-shrinking ranges \(\to\) HLD / Euler (\(\chref{c16},\chref{c24}\))
(\(\defref{c30-pick}\)).
\end{quickref}

\begin{lstlisting}[caption={Chapter 26 kit: the virtual-tree build (sort by tin, close under adjacent lca, dedup, monotone stack with depth-difference weights), the weighted fold, the pairwise distance sum, and a weighted-centroid descent -- one compact skeleton each, in one place. ASCII only, identifiers spelled without underscores.}, label={lst:c30-kit}]
// ===== BUILD the virtual tree (26.1) =====
// tin[], depth[] known; lca(u,v), isAncestor(u,v) from Ch.16
int vtSeen[MAXN]; vector<pair<int,int>> vtKids[MAXN]; // (child,L)
vector<int> buildVT(vector<int> S){
    sort(S.begin(),S.end(),[](int a,int b){return tin[a]<tin[b];});
    vector<int> N=S;
    for(int i=0;i+1<(int)S.size();i++) N.push_back(lca(S[i],S[i+1])); //k1
    sort(N.begin(),N.end(),[](int a,int b){return tin[a]<tin[b];});
    N.erase(unique(N.begin(),N.end()),N.end());         // |N|<=2k-1
    vector<int> st; st.push_back(N[0]);
    for(int i=1;i<(int)N.size();i++){ int v=N[i];
        while(!isAncestor(st.back(),v)) st.pop_back();
        int u=st.back(); addVEdge(u,v,depth[v]-depth[u]); // weight L
        st.push_back(v); }
    return N; }       // N[0] is the global lca

// ===== WEIGHTED FOLD (26.2) =====
ll cnt[MAXN], acc[MAXN];
void vtDfs(int v){ cnt[v]=vtSeen[v]; acc[v]=0;
    for(auto[u,L]:vtKids[v]){ vtDfs(u); cnt[v]+=cnt[u];
        acc[v]+=acc[u]+perEdge(cnt[u])*L; } }

// ===== PAIRWISE DISTANCE SUM (26.3) =====
ll pairSum(int v,int k){ ll res=0;
    for(auto[u,L]:vtKids[v])
        res += (ll)L*cnt[u]*(k-cnt[u]) + pairSum(u,k);
    return res; }

// ===== WEIGHTED CENTROID / 1-MEDIAN descent (Summit shape) =====
// walk from root into the child whose marked weight is a strict majority
int centroid(int root){ int v=root;
    while(true){ int best=-1,bw=0;
        for(auto[u,L]:vtKids[v]) if(cnt[u]*2>cnt[root] && cnt[u]>bw){bw=cnt[u];best=u;}
        if(best<0) return v; v=best; } }                 // then cost = sum w*dist(., centroid)
\end{lstlisting}

\chaptersection{Write-From-Memory Drills}

Five signature skeletons, each lifted from a listing already in the chapter. Cover the marked line,
reproduce from memory, then check the answer one line below.

\begin{lstlisting}[caption={Drill 1 -- close a marked set under lca (from lst:c30-build)}]
sort(S by tin);  vector<int> N = S;
for (int i=0; i+1 < (int)S.size(); i++)
    N.push_back( lca(S[i], S[i+1]) );     // k-1 ADJACENT pairs, not C(k,2)
sort(N by tin);
N.erase( unique(N.begin(), N.end()), N.end() );   // now N = A(S), |N| <= 2k-1
\end{lstlisting}
\noindent{\scriptsize Key: after sorting by \(\tin\), every \(\lca\) of marked nodes already appears
between two \emph{neighbours} in that order (\(\thmref{c30-size}\)); the second sort+unique removes the
duplicates a repeated LCA introduces (\(\coderef{c30-build}\)).}

\begin{lstlisting}[caption={Drill 2 -- the monotone-stack link (from lst:c30-build)}]
vector<int> st; st.push_back(N[0]);                    // N[0] = global lca = root
for (int i=1; i < (int)N.size(); i++){ int v = N[i];
    while (!isAncestor(st.back(), v)) st.pop_back();   // climb to nearest kept ancestor
    int u = st.back();
    vtKids[u].push_back({ v, depth[v] - depth[u] });    // compressed edge weight = L
    st.push_back(v);
}
\end{lstlisting}
\noindent{\scriptsize Key: the stack is the current kept root-path; popping until the top is an
ancestor finds \(v\)'s nearest ancestor in \(A(S)\) in amortised \(\bigO{1}\); the weight is the
\emph{depth difference}, the number of real edges the chain swallowed (\(\defref{c30-edgeweight}\)).}

\begin{lstlisting}[caption={Drill 3 -- the weighted fold (from lst:c30-fold)}]
void vtDfs(int v){
    cnt[v] = marked[v] ? 1 : 0;                         // count MARKED terminals
    acc[v] = 0;
    for (auto [u,L] : vtKids[v]){
        vtDfs(u);
        cnt[v] += cnt[u];
        acc[v] += acc[u] + perEdge(cnt[u]) * L;         // per-edge quantity x its real length
    }
}
\end{lstlisting}
\noindent{\scriptsize Key: identical to Ch.~13's subtree fold, except each child edge is multiplied by
its length \(L\) and \(\mathit{cnt}\) counts only marked terminals -- so a \(\bigO{n}\) DP runs in
\(\bigO{k}\) with the same answer (\(\thmref{c30-fold}\), \(\coderef{c30-fold}\)).}

\begin{lstlisting}[caption={Drill 4 -- pairwise distance sum over a set (from lst:c30-pairsum)}]
ll pairSum(int v, int k){ ll res = 0;
    for (auto [u,L] : vtKids[v])
        res += (ll)L * cnt[u] * (k - cnt[u])            // each of L edges splits cnt[u] | k-cnt[u]
               + pairSum(u, k);
    return res;
}
\end{lstlisting}
\noindent{\scriptsize Key: an edge is on a marked-to-marked path iff exactly one endpoint is below it;
\(L\) swallowed edges share the same split, so the term is \(L\cdot\mathit{cnt}\cdot(k-\mathit{cnt})\)
(\(\thmref{c30-pairsum}\), \(\coderef{c30-pairsum}\)).}

\begin{lstlisting}[caption={Drill 5 -- weighted-centroid / 1-median descent (from lst:c30-summit)}]
int centroid(int root){ int v = root;
    while (true){
        int best = -1, bw = 0;
        for (auto [u,L] : vtKids[v])
            if (cnt[u] * 2 > cnt[root] && cnt[u] > bw){ bw = cnt[u]; best = u; }
        if (best < 0) return v;                 // no strict-majority child -> we are the median
        v = best;
    }
}
\end{lstlisting}
\noindent{\scriptsize Key: stepping into a child that holds \(>\) half the total weight lowers the sum
of weighted distances; when no child does, the node is the weighted centroid minimising
\(\sum \mathit{weight}\cdot\dist\) -- the Summit's assembly node (\(\probref{c30-summit}\)).}

\chaptersection{Complete Problem Checklist}

Every problem this chapter offers -- two fully solved in-section, worked traces on the running
example \(T,S=\{4,8,6,10\}\), and one Summit on its own divisor-tree cases. Tick \emph{only} after an
unaided solve; a hinted attempt counts as half.

\begin{center}
\footnotesize
\setlength{\tabcolsep}{4pt}
\renewcommand{\arraystretch}{1.25}
\begin{tabular}{@{}c p{0.16\linewidth} l p{0.20\linewidth} >{\raggedright\arraybackslash}p{0.26\linewidth} c@{}}
\toprule
\textbf{\#} & \textbf{Problem} & \textbf{Source} & \textbf{Key technique} & \textbf{What to look for} & \textbf{Done?} \\
\midrule
\multicolumn{6}{l}{\textbf{Section 26.1 -- construction}} \\
W26.1 & stack sweep of \(A(\{4,8,6,10\})\) & worked & \lv{4} sort-by-tin + adjacent LCA & closed set \(1,2,4,8,3,6,10\); six weighted edges; \(\dist(8,10)=6\) & \(\square\) \\
W26.1b & Ancient Tree & CF 2127E & \lv{4} one VT per colour & branch node of \(2+\) colour VTs is forced cutie; star test \(\to\) cost \(0\) & \(\square\) \\
\midrule
\multicolumn{6}{l}{\textbf{Section 26.2 -- DP on the compressed tree}} \\
W26.2 & weighted fold equals the \(\bigO{n}\) DP & worked & \lv{5} edge weight + \(\mathit{cnt}\) & a length-\(L\) edge carries \(L\) copies of its flow; nodes off-paths dropped & \(\square\) \\
W26.2b & Spectral Components & CF 2252F & \lv{6} per-colour VT + count + greedy & virtual edge worth \(\min(m_{\text{in}},m_c-m_{\text{in}})\); test \(1\to\) \texttt{1 3 -1 -1 -1} & \(\square\) \\
\midrule
\multicolumn{6}{l}{\textbf{Section 26.3 -- applications (worked only)}} \\
W26.3 & pair-sum + Steiner size on \(S\) & worked & \lv{5} \(\sum w\cdot\mathit{cnt}(k-\mathit{cnt})\) & pair-sum \(=26\) over six pairs; Steiner \(=8\) edges & \(\square\) \\
\midrule
\multicolumn{6}{l}{\textbf{Section 26.4 -- recognition (worked only)}} \\
W26.4 & route five specs through the gate & worked & \lv{4} three-question signature & fixed / marked-induced / small \(\sum k\); reject graph Steiner & \(\square\) \\
\midrule
\multicolumn{6}{l}{\textbf{Section 26.5 -- the method boundary (worked only)}} \\
W26.5 & six specs a VT must hand away & worked & \lv{5} which engine & every-node \(\to\) DP; graph \(\to\) bitmask; moving \(\to\) LCT; range \(\to\) HLD & \(\square\) \\
\midrule
\multicolumn{6}{l}{\textbf{Summit (chapter end)}} \\
W26 & Chaotic V. & CF 1292D & \lv{6} divisor-tree VT + centroid & factorials closed by \(\tin\); descend majority branch; \(\to\) 5, 6, 6, 11 & \(\square\) \\
\bottomrule
\end{tabular}
\end{center}

\chaptersection{Recommended Study Order}

Four focused days. Day~1 internalises the build and CF 2127E; day~2 the weighted fold and CF
2252F; day~3 the two application families; day~4 recognition, the boundary, the Summit and the
mixed set.

\begin{center}
\small
\renewcommand{\arraystretch}{1.25}
\begin{tabular}{@{}llp{0.55\linewidth}@{}}
\toprule
\textbf{Day} & \textbf{Focus} & \textbf{Tasks} \\
\midrule
1 & Read + do 26.1 & hand-sweep the stack on \(S=\{4,8,6,10\}\) to the six weighted edges (\(\coderef{c30-build}\)); prove \(|A(S)|\le2k-1\) (\(\thmref{c30-size}\)); then CF 2127E, replaying the star test to cost \(0\) (\(\probref{c30-ancient}\)) \\
2 & Read + do 26.2 & write the weighted fold cold (\(\coderef{c30-fold}\)); then CF 2252F, replaying test 1 to \texttt{1 3 -1 -1 -1} (\(\probref{c30-spectral}\)) \\
3 & Read + do 26.3 & derive the pair-sum \(=26\) and Steiner \(=8\) by hand against the six direct distances (\(\thmref{c30-pairsum}\), \(\coderef{c30-pairsum}\)) \\
4 & Read + do 26.4, 26.5, Summit & place five specs on the gate (\(\defref{c30-signature}\)) and six on the pick (\(\defref{c30-pick}\)); then CF 1292D, replaying the descent to \(5,6,6,11\) (\(\probref{c30-summit}\)); finish with the first 3 mixed entries \\
\bottomrule
\end{tabular}
\end{center}

\chaptersection{Difficulty Progression}

Tiers are competency gates -- each row names the \emph{skill} the tier unlocks.

\begin{center}
\small
\renewcommand{\arraystretch}{1.3}
\begin{tabular}{@{}lp{0.32\linewidth}p{0.40\linewidth}@{}}
\toprule
\textbf{Tier} & \textbf{Problems} & \textbf{What passing this tier proves} \\
\midrule
\lv{3} \ \(\bigstar\bigstar\bigstar\)
  & close a marked set; name \(|A(S)|\)
  & you sort by \(\tin\) and take only the \(k-1\) adjacent LCAs, dedup included \\
\lv{4} \ \(\bigstar\bigstar\bigstar\bigstar\)
  & the stack build; CF 2127E (mixed at CF~2200)
  & you hang every node on its nearest kept ancestor with the right edge weight \\
\lv{5} \ \(\bigstar\bigstar\bigstar\bigstar\bigstar\)
  & weighted fold, pair-sum / Steiner, recognition, mixed at CF~2300--2500
  & you make a \(\bigO{k}\) virtual DP match the \(\bigO{n}\) honest one \\
\lv{6} \ \(\bigstar\bigstar\bigstar\bigstar\bigstar\bigstar\)
  & CF 2252F, CF 1292D (Summit), mixed at CF~2400--2700
  & you \emph{recognise} the marked-set signature (or reject it) and fold a weighted centroid \\
\bottomrule
\end{tabular}
\end{center}

\chaptersection{Reduction Gallery}

Every entry is ``one marked-set shape \(\to\) one virtual-tree move (or an honest hand-off).'' Each
line collapses a problem family onto a listing already in the chapter.

\begin{itemize}[nosep, leftmargin=1.4em]
  \item ``per colour / per query, an aggregate over the vertices of that set and the paths joining them''
        \(\to\) one virtual tree per set, Ch.~13 fold on it (\(\thmref{c30-fold}\), \(\coderef{c30-build}\)).
  \item ``is this vertex a forced branch point of colour \(d\)?'' (CF 2127E) \(\to\) it is an internal node
        of colour \(d\)'s virtual tree; count memberships (\(\probref{c30-ancient}\)).
  \item ``sum of distances from a set to a chosen connected component'' (CF 2252F) \(\to\) per-colour VT,
        edge worth \(\min(m_{\text{in}},m_c-m_{\text{in}})\), greedy \(k_c-1\) heaviest (\(\probref{c30-spectral}\)).
  \item ``sum of \(\dist\) over all pairs of a \(k\)-set'' \(\to\) per-edge separation \(\sum w\cdot\mathit{cnt}(k-\mathit{cnt})\)
        (\(\thmref{c30-pairsum}\), \(\coderef{c30-pairsum}\)).
  \item ``minimum connected subgraph spanning \(S\) \emph{on a tree}'' \(\to\) the expanded virtual tree,
        size \(\sum\mathit{w}(e)\) -- forced, no search (\(\secref{c30-apps}\)).
  \item ``place a facility minimising the weighted sum of distances to the marked nodes'' (CF 1292D)
        \(\to\) weighted centroid / 1-median on \(A(S)\) (\(\probref{c30-summit}\), \(\coderef{c30-summit}\)).
  \item ``the value is summed over \emph{every} node'' \(\to\) \emph{not} a virtual tree: full tree DP
        or rerooting (\(\defref{c30-pick}\), \(\secref{c30-boundary}\)).
  \item ``connect \(k\) terminals on a \emph{graph}'' \(\to\) \emph{not} a virtual tree: they meet
        at new nodes, so bitmask Steiner DP (\(\secref{c30-boundary}\), mixed CF 152E).
  \item ``path/subtree ranges on a fixed tree, nothing shrinks with \(k\)'' \(\to\) HLD / Euler+Fenwick
        (\(\chref{c16}\), \(\chref{c24}\)); ``edges are edited online'' \(\to\) link-cut (\(\chref{c23}\)).
\end{itemize}

\chaptersection{Mixed Problem Set}

Ten entries, each live-verified against the current Codeforces problem set (names, ratings and tags
re-pulled this session) and absent from every other roster in the book (the famous virtual-tree problems
CF 613D in Ch.~27 and CF 809E in Ch.~25 are owned elsewhere and not re-linked here). None says ``use a
virtual tree''; several deliberately \emph{do not} want one, and the drill is to run them through the
signature (\(\defref{c30-signature}\)) and say so. Each hint names the reduction (or the rejection) and
the Chapter~26 listing it leans on, ordered ascending by rating.

\begin{mixedset}
  \practice{Compressed Tree}{CF 1901E}
    {https://codeforces.com/problemset/problem/1901/E}{CF 2200 L5}
    {Tags \emph{dfs, dp, greedy, trees}. The name and shape -- a big tree, a small set of special
     vertices, an answer that routes through the paths joining them -- is the signature firing on all
     three questions: fixed tree, marked-set-induced answer, \(\sum k\) small (\(\defref{c30-signature}\)).
     Compress to \(A(S)\) and run the Ch.~13 fold on it (\(\coderef{c30-build}\), \(\thmref{c30-fold}\)).
     This is the genuine virtual-tree case of the set.}
  \practice{Diameter Intersections}{CF 2253E}
    {https://codeforces.com/problemset/problem/2253/E}{CF 2200 L5}
    {Tags \emph{brute force, dfs, fft, graphs, trees}. The answer is about the tree's \emph{diameter}
     and paths crossing it -- a global structural quantity, not a per-set fold. Gate S2 fails: the object
     is defined on all nodes, so nothing marked-set-induced shrinks (\(\defref{c30-pick}\)). Use this to
     practise REJECTING the virtual tree; the real tool is a diameter analysis (the \emph{fft} tag is a
     red herring for the hard sub-case). Contrast 26.5's first spec.}
  \practice{Tree Colorings}{CF 2112E}
    {https://codeforces.com/problemset/problem/2112/E}{CF 2200 L5}
    {Tags \emph{combinatorics, dfs, dp, math, number theory}. A counting problem on colourings with a
     combinatorics/number-theory core; the DP is over the whole tree's structure, not over a small
     marked subset's induced subtree. Ask the gate: does dropping the unmarked nodes lose information?
     If the count depends on every node's colour choice, it is a plain tree-DP, not a virtual tree
     (\(\defref{c30-signature}\), \(\secref{c30-boundary}\)).}
  \practice{Nightcrawler}{CF 2254G}
    {https://codeforces.com/problemset/problem/2254/G}{CF 2200 L5}
    {Tags \emph{data structures, dfs, dp, greedy, sortings}. A greedy/sorting tree problem with a DFS
     order -- a strong hint that \(\tin\) matters, but check whether the answer is confined to a small
     marked set (VT) or is a global greedy over all nodes (tree DP). If each query names a tiny set whose
     induced subtree you re-fold, it collapses to \(\bigO{k\log k}\) via \(\coderef{c30-build}\); if the
     greedy runs over the whole tree, it does not shrink and stays a Ch.~13 job (\(\defref{c30-pick}\)).}
  \practice{Build a Tree}{CF 2247E}
    {https://codeforces.com/problemset/problem/2247/E}{CF 2300 L5}
    {Tags \emph{constructive algorithms, trees, two pointers}. A constructive task: assign values so
     constraints hold. There is no marked set whose aggregate you compute -- you \emph{build} an
     assignment. Recognition reflex: a construction is rarely a virtual tree; solve by the greedy / two
     pointer the tags advertise. A deliberate rejection (\(\secref{c30-boundary}\)'s warning bullet 1:
     named vertices, but no set-induced sum to compress).}
  \practice{Dogeforces}{CF 1494D}
    {https://codeforces.com/problemset/problem/1494/D}{CF 2300 L5}
    {Tags \emph{constructive, dsu, divide and conquer, dfs}. You reconstruct a hidden tree from a
     distance-like matrix with DSU and divide-and-conquer -- the whole node set is the object, not a
     marked subset, and the tree is being \emph{built} not given-and-compressed. Reject the virtual tree
     (\(\defref{c30-signature}\): S1 has no fixed tree yet); reach for DSU + the reconstruction.}
  \practice{Again Trees\ldots (Easy Version)}{CF 2194F1}
    {https://codeforces.com/problemset/problem/2194/F1}{CF 2300 L6}
    {Tags \emph{bitmasks, dfs, dp, trees}. The \emph{bitmasks} tag is the tell: this is a
     subset-DP, exactly the shape 26.5 routes to bitmask DP rather than a virtual tree. On a \emph{tree}
     a small marked set would compress free (\(\secref{c30-apps}\)); when the state must range over
     \(2^k\) terminal subsets the answer does \emph{not} live on a single induced subtree, so run the
     bitmask fold. Distinguish it from CF 2127E where the set \emph{is} fixed per colour.}
  \practice{Spanning Tree Queries}{CF 1633E}
    {https://codeforces.com/problemset/problem/1633/E}{CF 2400 L5}
    {Tags \emph{binary search, dsu, graphs, greedy, sortings}. A spanning-forest / bottleneck question
     answered by sorting edges, DSU, and binary-searching the threshold -- a \emph{graph} MST-shape, not
     a fixed-tree marked-set fold. Gate S1 fails (the structure is a graph and you choose edges), so the
     virtual tree is the wrong lens; hand it to DSU + binary search (\(\defref{c30-pick}\)).}
  \practice{Garden}{CF 152E}
    {https://codeforces.com/problemset/problem/152/E}{CF 2500 L5}
    {Tags \emph{bitmasks, dp, graphs, trees}. The canonical Steiner-on-a-graph problem: connect a small
     set of terminals in a weighted \emph{graph}, and the cheapest connector can be a non-terminal. That
     is precisely the ``graph, new branch nodes'' rejection of \(\secref{c30-boundary}\) -- solve with the
     Dreyfus--Wagner bitmask Steiner DP \(O(3^k n)\), \emph{not} a virtual tree. On a plain \emph{tree} the
     same question would be free via the expanded \(A(S)\) (\(\secref{c30-apps}\)); the graph is the veto.}
  \practice{Tree Queries}{CF 1254D}
    {https://codeforces.com/problemset/problem/1254/D}{CF 2700 L5}
    {Tags \emph{data structures, probabilities, trees}. A probabilistic tree question solved with a
     data-structure sweep over a \emph{fixed} full tree -- the marked set is not the driver of the answer,
     and expectation runs over global structure. Nothing here is per-subset compression; treat it as the
     sharpest rejection at the top of the band (compare \(\secref{c30-boundary}\): a big tree and an
     aggregate over the whole thing, so \(A(S)\) would drop mass the expectation needs).}
\end{mixedset}



\chaptersection{The Summit -- CF 1292D Chaotic V.}

\begin{notebox}[Why this is the chapter top rope]
Every in-section problem here compressed a marked set on a tree you were \emph{given} as
\(n\) explicit nodes. The Summit hands you a tree you can never write down: the nodes
are all positive integers, \(x>1\) joins \(x/\!\text{lpf}(x)\) (lowest prime divisor), and
the fragments sit at factorial nodes \(k_i!\) with \(k_i\le 5000\) -- so the tree is
astronomically deep (\(\Omega(5000!)\approx 30{,}000\)) and its node labels overflow any
machine integer. The compression is the whole solution: only the \(\le 5000\) \emph{distinct}
factorials matter, their \(\lca\) are the common \emph{descending-prime-factor prefixes}, so
the virtual tree is that trie, and the answer -- minimise the sum of path lengths to an
assembly node \(P\) -- is a \textbf{weighted centroid / 1-median} on it
(\(\thmref{c30-fold}\) for the fold, \(\secref{c30-apps}\) for the distance currency, and the
Ch.~24 centroid idea for the descent). It is the honest summit, not the highest badge: the
hardest problem this chapter fully derives and hand-replays to all four judge outputs
(\(5,6,6,11\)).
\end{notebox}

\solvedproblem{c30-summit}{Chaotic V.}{CF 1292D}
  {https://codeforces.com/problemset/problem/1292/D}{CF 2700 L6}
  {implicit divisor tree (node \(x\!>\!1\) attaches to \(x/\!\text{lpf}(x)\), root \(1\),
   depth \(\depth(x)=\Omega(x)\)); the marked set is the distinct factorials \(k_i!\) with
   multiplicities; their LCAs are common prime-factor prefixes in DECREASING order, so the
   virtual tree is that trie; the objective is the WEIGHTED CENTROID / 1-median, reached by
   Akikaze's descent -- repeatedly step into the branch holding a strict majority of fragments,
   tracking the cost delta \(\mathit{cur}-=\mathit{majority}\), \(\mathit{cur}+= (n-\mathit{majority})\)}

\begin{problemstatement}[CF 1292D Chaotic V. (statement, constraints and all four samples verified against the live CF page)]
The A.R.C.\ Library core is an undirected graph of \emph{infinitely many} nodes labelled by
all positive integers \(1,2,3,\dots\). The node \(x\) (\(x>1\)) is directly connected to the
node \(x/f(x)\), where \(f(x)\) is the \textbf{lowest prime divisor} of \(x\). (This makes a
tree rooted at \(1\): repeatedly dividing by the lowest prime factor peels primes in
increasing order and reaches \(1\), so \(\depth(x)=\Omega(x)\), the number of prime factors of
\(x\) with multiplicity.)

Vanessa's mind is split into \(n\) fragments; the \(i\)-th sits at the node \(k_i!\) (a
factorial). Choose a node \(P\) (the assembly node) so that the \emph{total length of the
paths from every fragment to \(P\)} is as small as possible; if several fragments share a
node, that node's distance to \(P\) is counted once per fragment. Output that minimum total.

\textbf{Input.} Line~1: \(n\) (\(1\le n\le 10^6\)). Line~2: \(n\) integers \(k_1,\dots,k_n\)
(\(0\le k_i\le 5000\)); fragment \(i\) is at \(k_i!\). Note \(0!=1!=1\).

\textbf{Output.} One integer -- the minimum sum of path lengths to \(P\).

\textbf{Samples (all four official).} \(n{=}3,\ k{=}[2,1,4]\to \mathbf{5}\).
\(n{=}4,\ k{=}[3,1,4,4]\to\mathbf{6}\). \(n{=}4,\ k{=}[3,1,4,1]\to\mathbf{6}\).
\(n{=}5,\ k{=}[3,1,4,1,5]\to\mathbf{11}\). (The note draws the first \(24\) nodes and states
the first answer uses \(P=1\), the second uses \(P=6=3!\).)
\end{problemstatement}

\begin{attemptgate}[Attempt it yourself first -- \lv{6}, budget a real sitting]
The tree is infinite and the fragment nodes are factorials up to \(5000!\). You cannot build
it, cannot even label its nodes. What small set of nodes could the answer possibly depend on?
\peek{an optimal \(P\) lies on at least one fragment-to-root path, so it lies in the union of
the paths between the fragment nodes -- exactly the virtual tree of the marked set
\(\{k_i!\}\). How many \emph{distinct} factorials can there be, given \(k_i\le 5000\)?}
\nudge{collapse the \(n\) fragments to counts \(\mathit{cnt}[k]\) over the \(\le 5001\) distinct
factorials. The objective \(\sum_v \mathit{weight}(v)\cdot\dist(v,P)\) is minimised at the
\textbf{weighted centroid} of the virtual tree: from any node, moving to a child that holds
\(>\tfrac{n}{2}\) of the total weight strictly lowers the sum, so walk into the majority child
until none is. Distances are \(\Omega\)-differences: \(\dist(x,y)=\Omega(x)+\Omega(y)-2\Omega(\lca)\).}
\fullhint{Two nodes branch by their prime factors in DECREASING order: they share the branch at
depth \(d\) iff their \(d{+}1\) largest prime factors agree. Hence the virtual tree of the
factorials is a \emph{trie} of each \(k_i!\)'s prime factors sorted largest-first, and its edges
have length \(1\) (one prime peeled). Precompute \(\mathit{expo}[i][p]\) = exponent of prime \(p\)
in \(i!\) cumulatively (\(\bigO{K^2}\), \(K=5000\)); then \(\Omega(i!)=\sum_p\mathit{expo}[i][p]\).
Start at \(P=1\) with \(\mathit{cur}=\sum_i \mathit{cnt}[i]\cdot\Omega(i!)\). Let each active
\(i\)'s \emph{current largest remaining prime} be its branch key; group fragment counts by that
key, step \(P\) down into the prime \(p\) holding a strict majority (\(\mathit{best}\cdot 2>n\)),
updating \(\mathit{cur}\mathrel{-}=\mathit{best}\), \(\mathit{cur}\mathrel{+}= (n-\mathit{best})\),
and \(\mathit{ans}=\min(\mathit{ans},\mathit{cur})\); then peel one \(p\) off every node that took
that branch and discard the nodes that did not. Stop when no branch is a majority -- \(\mathit{ans}\)
is the 1-median cost. Read \(\mathit{cur}\) as \code{long long}.}
\end{attemptgate}

\solstep{Understanding}
The infinite divisor tree is a tree the problem never lets you touch: labels blow past any
integer type and depth reaches \(30{,}000\). The rescue is that the answer depends on only the
marked nodes -- the fragments at \(k_i!\) -- and the paths joining them, i.e.\ their virtual
tree (\(\defref{c30-vt}\)). Two facts make that tree explicit and small. First, collapsing
identical factorials to a weight \(\mathit{cnt}[k]\) leaves \(\le 5001\) marked nodes.
Second, the ancestor relation is pure prime-factor bookkeeping: node \(x\)'s path to \(1\) is
read off the prime factors of \(x\) in \emph{decreasing} order (peel the largest first when you
descend from the root), so \(x\) is an ancestor of \(y\) iff \(x\)'s decreasing-prime list is a
prefix of \(y\)'s, and \(\lca(x,y)\) is their common prefix. Thus the virtual tree of the
factorials is the \term{trie} of their decreasing-prime-factor sequences, every edge of length
\(1\), with a weight \(\mathit{cnt}\) parked at the factorial leaves.

The objective is then a textbook 1-median: \(\sum_v \mathit{weight}(v)\cdot\dist(v,P)\) over tree
nodes is minimised at the \textbf{weighted centroid} -- the unique node after which no child
subtree holds a strict majority of the total weight -- because shifting \(P\) across an edge into
a side holding \(t\) of \(n\) units changes the sum by \((n-t)-t=n-2t\), negative exactly when
\(t>n/2\). So the whole problem is ``build the (implicit) virtual tree, then walk to its weighted
centroid, tracking the cost.'' That is \(\thmref{c30-fold}\)'s principle (a \(\bigO{k}\) structure
replaces the \(\bigO{n}\) tree) fused with Chapter~24's centroid move.

\solstep{Key Observations}
\begin{observation}
\textbf{Depth is a prime count.} \(\depth(x)=\Omega(x)\) (prime factors with multiplicity), so
\(\depth(k!)=\Omega(k!)=\sum_{j=2}^{k}\Omega(j)\), precomputed cumulatively in \(\bigO{K^2}\) with a
smallest-prime sieve (\(\defref{c30-edgeweight}\)'s currency, here unit-length edges).
\end{observation}
\begin{observation}
\textbf{Factorials are NOT a chain.} \(4!=24\) is not an ancestor of \(5!=120\): the root-to-\(120\)
path is \(1,5,15,30,60,120\) (primes \(5,3,2,2,2\)) and never visits \(24\). The induced structure
is the trie of \emph{decreasing} prime-factor sequences, so you must group by common prefix, not
assume \(k!\mid (k{+}1)!\) gives tree ancestry. (\(\thmref{c30-size}\) still caps the virtual tree
at \(2k-1\) marked-plus-LCA nodes.)
\end{observation}
\begin{observation}
\textbf{The optimum is the weighted centroid.} For the sum of weighted distances the move rule is
local: cross an edge into a component with \(>n/2\) weight and the total drops by \(2t-n>0\).
Since \(P\) must lie in the marked Steiner skeleton, the descent is confined to the virtual tree,
and the accumulated \(\mathit{ans}\) over the walk is the minimum (\(\probref{c30-summit}\)).
\end{observation}

\solstep{Think It With Me}
\begin{thinkbox}
\deadend{Materialise the divisor tree up to \(\max k_i!\) and DFS for the centroid}
{The nodes are labelled by integers as large as \(5000!\) -- they do not fit in \code{long long},
and the depth is \(\Omega(5000!)\approx 3\times10^4\), so the tree has far more nodes than you can
store or reach. You can never write the tree down; only the \(\le5001\) marked factorials and
their prefix-LCAs are addressable. \textbf{TLE / overflow}.}
\deadend{Assume the factorials form a chain \(1!\mid2!\mid\cdots\) and run a line 1-median}
{\(k!\) divides \((k{+}1)!\) as integers, but divisibility is NOT ancestry in this tree: \(4!=24\)
is not on the path to \(5!=120\). Treating the marked set as a line puts the median at the wrong
node and over-counts every branch that is really a shared prime prefix. This is the seductive WA:
the answer is a \emph{centroid on a trie}, not a median on a path. \textbf{WA}.}
\deadend{Recompute each factorial's full factorisation and its LCA per pair, on the fly, all \(n\)}
{With \(n\le10^6\) fragments and \(\Omega\) up to thousands, doing \(O(\Omega)\) work per fragment
per step is \(\tle\). You must (i) collapse to \(\mathit{cnt}[k]\) over \(\le5001\) distinct
\(k\), and (ii) descend the majority branch by grouping \emph{distinct} factorials, never the
\(10^6\) raw fragments. \textbf{TLE}.}
\aha{Collapse to the \(\le5001\) distinct factorials with multiplicities, note the virtual tree is
the trie of their \emph{decreasing} prime-factor sequences (unit edges), and walk the weighted
centroid: from \(P=1\) with \(\mathit{cur}=\sum_k\mathit{cnt}[k]\,\Omega(k!)\), repeatedly step into
the branch whose current-largest-prime group is a strict majority, updating
\(\mathit{cur}\mathrel{-}=\mathit{best}\) and \(\mathit{cur}\mathrel{+}=n-\mathit{best}\), record
\(\min\), and peel one prime off the nodes that took that branch. The running minimum is the 1-median
cost.}
\end{thinkbox}

\solstep{The Approach Ladder}
\begin{center}\renewcommand{\arraystretch}{1.25}
\begin{tabular}{@{}c p{0.23\linewidth} p{0.29\linewidth} l c@{}}
\toprule
\textbf{\#} & \textbf{Approach} & \textbf{Idea} & \textbf{Time} & \textbf{Verdict}\\
\midrule
1 & \lv{2}\ build the tree to \(\max k_i!\), DFS & write every node & -- & \textcolor{warningred}{\ding{55}\ \textbf{RE}} \\
2 & \lv{3}\ line 1-median over \(1!,2!,\dots\) & assume factorials are a chain & \(\bigO{n}\) & \textcolor{warningred}{\ding{55}\ \textbf{WA}} \\
3 & \lv{4}\ pairwise-LCA virtual tree, raw \(n\) & right model, no collapsing & \(\bigO{n\Omega}\) & \tle \\
4 & \lv{6}\ \textbf{trie of decreasing primes + centroid descent} & collapse to \(\le K\) counts, walk majority & \(\bigO{K^2 + n}\) & \ac \\
\bottomrule
\end{tabular}\end{center}
Rung~4 is the Summit: rungs~1--3 have the right \emph{goal} but the wrong \emph{object} (an
impossible tree, a false chain, or an un-collapsed set); realising the virtual tree is a
\emph{decreasing-prime-factor trie} and that the optimum is its weighted centroid is the whole win.

\solstep{Solution}
Read \(n\) and collapse the fragments into counts \(\mathit{cnt}[k]\) for \(0\le k\le 5000\).
Precompute \(\mathit{expo}[i][p]\) = the exponent of prime \(p\) in \(i!\), by carrying the row
of \((i{-}1)!\) forward and adding the factorisation of \(i\) (a \(\bigO{K^2}\) sieve-style DP).
Initialise \(P\) at the root: \(\mathit{cur}=\mathit{ans}=\sum_i \mathit{cnt}[i]\sum_p\mathit{expo}[i][p]\)
(that is \(\sum_i\mathit{cnt}[i]\,\Omega(i!)\)), and set each active \(i\)'s pointer
\(\mathit{bestPD}[i]\) to the largest prime still present in \(i!\). Now descend: while some
pointer exceeds \(1\), group all fragment counts by their pointer prime, take the prime
\(\mathit{bestPrime}\) with the largest group \(\mathit{bestGroup}\); if \(\mathit{bestGroup}\cdot2\le n\)
or \(\mathit{bestPrime}=1\), stop (no strict majority). Otherwise move \(P\) one edge down into
that branch: \(\mathit{cur}=\mathit{cur}-\mathit{bestGroup}+(n-\mathit{bestGroup})\),
\(\mathit{ans}=\min(\mathit{ans},\mathit{cur})\); then every node whose pointer equalled
\(\mathit{bestPrime}\) drops one copy of that prime and reveals its next-largest, and every node
that did NOT take the branch is parked (\(\mathit{bestPD}=1\)). Print \(\mathit{ans}\).

\solstep{Proof of Correctness}
\begin{proofbox}
\textbf{(Virtual tree is the right object.)} An assembly node minimising the sum of distances lies
in the union of paths between marked nodes (otherwise slide it toward that union and every term
drops), so confining the search to the virtual tree of \(S=\{k_i!\}\) loses no optimum
(\(\thmref{c30-fold}\)). \textbf{(Structure of that tree.)} In the divisor tree the root-to-\(x\)
path is the sequence of partial products of \(x\)'s prime factors in decreasing order, hence
\(x\) is an ancestor of \(y\) iff \(x\)'s decreasing list is a prefix of \(y\)'s; so the virtual
tree of any node set is the trie of their decreasing-prime-factor lists with unit edges, and
\(\depth(x)=\Omega(x)\) makes the root-cost \(\sum_i\mathit{cnt}[i]\Omega(k_i!)\) exact.
\textbf{(Descent finds the centroid.)} Crossing an edge into the side with \(t\) of \(n\) weight
changes the sum by \(n-2t\), which is negative iff \(t>n/2\); the loop moves exactly when a
strict-majority branch exists and takes the (necessarily unique) such branch, so each step
strictly lowers the sum and no local step is skipped -- the walk halts at the unique weighted
centroid, the global minimiser. \textbf{(Cost tracking.)} \(\mathit{cur}\) is updated by the same
\(-t+(n-t)\) rule at each real edge, and \(\mathit{ans}\) records the running minimum, which is
the value at the final centroid. \textbf{(Grouping by largest prime = descending the trie.)}
Partitioning active nodes by \(\mathit{bestPD}\) and peeling the majority prime advances every
relevant node down one trie level and freezes the others, maintaining the invariant that
\(\mathit{bestPD}[i]\) is the next trie edge below \(P\) on \(i!\)'s path. Hence \(\mathit{ans}\)
equals the true minimum sum of path lengths. \(\Box\)
\end{proofbox}

\solstep{Implementation}
\begin{lstlisting}[caption={CF 1292D Summit: collapse fragments to counts over the distinct factorials, precompute prime exponents of each factorial, then descend the decreasing-prime-factor trie into the strict-majority branch (weighted centroid) tracking the cost delta. Faithful to Akikaze's official solution. ASCII only, identifiers without underscores.}, label={lst:c30-summit}]
#include <bits/stdc++.h>
using namespace std;
typedef long long ll;
const int MAXK = 5000;
int n, cnt[MAXK+1];
int expo[MAXK+1][MAXK+1];            // expo[i][p] = exponent of prime p in i!
int main(){
    scanf("%d", &n);
    for (int i=0;i<n;i++){ int k; scanf("%d",&k); cnt[k]++; }
    // cumulative prime exponents of i!
    for (int i=2;i<=MAXK;i++){
        for (int p=0;p<=MAXK;p++) expo[i][p] += expo[i-1][p];
        int t=i;
        for (int x=2; (ll)x*x<=t; x++) while(t%x==0){ expo[i][x]++; t/=x; }
        if (t>1) expo[i][t]++;
    }
    // P at root: cost = sum over fragments of Omega(i!) = sum of exponents
    ll ans=0, cur=0; int bestPD[MAXK+1];
    for (int i=0;i<=MAXK;i++) bestPD[i]=1;
    for (int i=1;i<=MAXK;i++){
        if (cnt[i]==0){ bestPD[i]=1; continue; }
        for (int p=1;p<=MAXK;p++){
            ans += (ll)expo[i][p]*cnt[i];
            cur += (ll)expo[i][p]*cnt[i];
            if (expo[i][p]!=0) bestPD[i]=p;      // largest prime still in i!
        }
    }
    // descend the trie into the strict-majority branch (weighted centroid)
    while (true){
        int mx=1; for (int i=0;i<=MAXK;i++) if (bestPD[i]>mx) mx=bestPD[i];
        if (mx<=1) break;
        int freq[MAXK+1]; memset(freq,0,sizeof freq);
        for (int i=0;i<=MAXK;i++) freq[bestPD[i]] += cnt[i];
        int bestPrime=1, bestGroup=freq[1];
        for (int p=2;p<=MAXK;p++) if (freq[p]>bestGroup){ bestGroup=freq[p]; bestPrime=p; }
        if (bestGroup*2 <= n) break;             // no branch holds a strict majority
        if (bestPrime==1) break;                 // majority is the parked/root group
        cur -= bestGroup; cur += (n - bestGroup);
        if (cur < ans) ans = cur;
        for (int i=1;i<=MAXK;i++){               // peel one bestPrime off the branch
            if (bestPD[i]!=bestPrime) bestPD[i]=1;      // others leave the descent
            if (bestPD[i]==1) continue;
            expo[i][bestPD[i]]--;
            while (bestPD[i]>1 && expo[i][bestPD[i]]==0) bestPD[i]--;
        }
    }
    printf("%lld\n", ans);
    return 0;
}
// Only the <=5001 distinct factorials are ever touched; the infinite tree is the trie of their
// prime factors in DECREASING order. Each descent step is one edge of that virtual tree.
\end{lstlisting}

\solstep{Trace on the Official Samples}
\begin{tracebox}[Sample 2: collapse, then descend to the centroid \(P=6\)]
Weights: \(\mathit{cnt}[3]{=}1\) (at \(3!{=}6\)), \(\mathit{cnt}[1]{=}1\) (at \(1!{=}1\)),
\(\mathit{cnt}[4]{=}2\) (at \(4!{=}24\)); \(n{=}4\), half \(=2\).
Depths \(\Omega\): \(\Omega(1!){=}0,\ \Omega(3!){=}2,\ \Omega(4!){=}4\). Root cost
\(\mathit{cur}=1{\cdot}2 + 1{\cdot}0 + 2{\cdot}4 = 10\), \(\mathit{ans}{=}10\).
Largest primes: \(3!{=}2{\cdot}3\to3\); \(4!{=}2^3{\cdot}3\to3\); \(1!\to\) parked.
\textbf{Step 1:} group by pointer: prime \(3\) holds \(\mathit{cnt}[3]{+}\mathit{cnt}[4]=3>n/2\).
Move into it: \(\mathit{cur}=10-3+(4-3)=8\), \(\mathit{ans}=8\) (\(P=3\)). Peel one \(3\):
\(3!\to2\) (\(\Omega\) left \(1\)), \(4!\to2\). Park node \(1\) (\(\mathit{bestPD}[1]{=}1\)).
\textbf{Step 2:} group by pointer: prime \(2\) holds \(\mathit{cnt}[3]{+}\mathit{cnt}[4]=3>n/2\).
Move into it: \(\mathit{cur}=8-3+1=6\), \(\mathit{ans}=6\) (\(P=6\)). Peel one \(2\): \(3!\to1\)
(parked), \(4!\to2\) (now largest prime \(2\) with one fewer exponent, still descending).
\textbf{Step 3:} group by pointer: prime \(2\) holds only \(\mathit{cnt}[4]=2\); \(2\cdot2\not>n\)
\(\Rightarrow\) \textbf{stop}. Assembly node \(P=6\), answer \(\boxed{6}\) -- matching the note
(\(P=6\), distances \(0,2,2,2\)). \(\Box\)
\end{tracebox}

\begin{tracebox}[The other three, by the same root-cost / majority rule]
\textbf{Sample 1} \(([2,1,4])\): \(\mathit{cnt}\) at \(2!{\to}1,1!{\to}1,4!{\to}1\); root cost
\(\Omega(2!){+}\Omega(1!){+}\Omega(4!)=1+0+4=5\). Largest primes \(2,(-),3\) tie at group size
\(1\), never a strict majority of \(n{=}3\) (need \(>1.5\Rightarrow2\)) \(\Rightarrow\) stop at
\(P=1\), answer \(\boxed{5}\).\\
\textbf{Sample 3} \(([3,1,4,1])\): \(\mathit{cnt}[3]{=}1,\mathit{cnt}[1]{=}2,\mathit{cnt}[4]{=}1\);
root cost \(2+0+4=6\). Prime \(3\) group \(=\mathit{cnt}[3]{+}\mathit{cnt}[4]=2\), not \(>2\)
\(\Rightarrow\) stop at \(P=1\), answer \(\boxed{6}\).\\
\textbf{Sample 4} \(([3,1,4,1,5])\): \(\mathit{cnt}[1]{=}2,\mathit{cnt}[3]{=}1,\mathit{cnt}[4]{=}1,\)
\(\mathit{cnt}[5]{=}1\); root cost \(0{+}2{+}4{+}5=11\). Prime groups: \(3\to\{3,4\}=2\),
\(5\to\{5\}=1\); max \(2\), not \(>2.5\) \(\Rightarrow\) stop at \(P=1\), answer \(\boxed{11}\).
All four match the judge. \(\Box\)
\end{tracebox}

\solstep{Complexity}
\begin{complexitybox}
Reading and collapsing to counts is \(\bigO{n}\). The prime-exponent DP over \(i!\) for
\(i\le K=5000\) is \(\bigO{K^2}\) time and \(\bigO{K^2}\) integers of table (\(\approx2.5\cdot10^7\),
about \(100\) MB, inside the 512 MB limit). The centroid descent runs at most one step per peeled
prime along the majority path -- \(\bigO{K^2}\) in the worst case but tiny in practice because it
halts the moment no branch is a strict majority -- with \(\bigO{K}\) work per step. All sums fit in
\code{long long} (\(n\le10^6\), \(\Omega(5000!)\approx3\times10^4\), product \(\le3\times10^{10}\)).
Total \(\bigO{n+K^2}\). \textbf{Space} \(\bigO{K^2}\). \emph{No node of the real tree is ever
built} -- the compression to the \(\le K\) distinct factorials' decreasing-prime trie is the entire
trick (\(\thmref{c30-fold}\)).
\end{complexitybox}

\begin{hardeningbox}
\begin{itemize}[nosep,leftmargin=1.4em]
  \item \textbf{Do not treat factorials as a chain.} \(4!\nmid_{\text{tree}}5!\) (24 is not on the
        path to 120); the ancestry is by \emph{decreasing-prime prefix}. Line-median logic is the
        Summit's signature \textbf{WA} (\(\probref{c30-summit}\)'s second dead end).
  \item \textbf{Park the non-majority nodes, and stop on ties.} A node that did not take the
        majority prime leaves the descent (\(\mathit{bestPD}=1\)); and the walk must halt when
        \(\mathit{bestGroup}\cdot 2\le n\) -- a group of \emph{exactly} half is not a strict
        majority, so moving would not lower the sum.
  \item \textbf{Collapse before you compute.} \(\mathit{cnt}[k]\) over \(\le5001\) distinct \(k\),
        never the \(10^6\) raw fragments; per-fragment factorisation is \tle (third dead end).
  \item \textbf{Use \code{long long} for the cost, int for exponents.} The running sum
        \(\mathit{cur}\) is up to \(n\cdot\Omega(K!)\approx3\times10^{10}\); only the exponent
        table fits in \code{int}. Overflow here silently corrupts the centroid comparison.
\end{itemize}
\end{hardeningbox}

\begin{recallbox}[Chapter 26 in ten seconds]
A virtual tree closes a marked set \(S\) under \(\lca\) with \(k-1\) adjacent-pair calls after a
\(\tin\)-sort, dedups to \(|A(S)|\le2k-1\), and links it with one monotone stack giving each edge
weight \(\depth(v)-\depth(u)\) (\(\thmref{c30-size}\), \(\coderef{c30-build}\)). Run Ch.~13's fold
on that compressed tree and you get the honest \(\bigO{n}\) answer in \(\bigO{k}\)
(\(\thmref{c30-fold}\), \(\coderef{c30-fold}\)); pairwise sums are \(\sum w\cdot\mathit{cnt}(k-\mathit{cnt})\)
and tree-Steiner is the expanded \(A(S)\) (\(\thmref{c30-pairsum}\), \(\coderef{c30-pairsum}\)).
Fixed-tree + marked-set-induced + small \(\sum k\) is the gate (\(\defref{c30-signature}\)); else
hand off to tree DP, bitmask Steiner, HLD or link-cut (\(\defref{c30-pick}\)). Summit CF 1292D
(\(\probref{c30-summit}\)) compresses an infinite divisor tree to a decreasing-prime-factor trie and
reads its weighted centroid.
\selftest{
  \item Why is sorting by \(\tin\) and adding only adjacent-pair LCAs enough to close \(S\)? \hfill\lv{4}
  \item What weight does a compressed edge \((u,v)\) carry, and what does \(\mathit{cnt}(v)\) count? \hfill\lv{5}
  \item State the pairwise-distance-sum formula and the running-example value. \hfill\lv{5}
  \item Give the two rejected shapes a virtual tree must hand away, and to what. \hfill\lv{5}
  \item In the Summit, why are the factorials a \emph{trie}, not a \emph{chain}, and where is the optimum? \hfill\lv{6}
}
\end{recallbox}

\chaptersection{Tool Kit Inventory}

The machine assembled through Chapter~26. Earlier chapters' rows are carried forward; this
chapter's are appended. Pointers only.

\begin{center}
\small
\renewcommand{\arraystretch}{1.25}
\begin{tabular}{@{}p{0.24\linewidth} l c p{0.22\linewidth} >{\raggedright\arraybackslash\scriptsize}p{0.17\linewidth}@{}}
\toprule
\textbf{Tool} & \textbf{Since} & \textbf{Lv} & \textbf{Trigger phrase} & \textbf{Where} \\
\midrule
\multicolumn{5}{l}{\itshape Carried forward} \\
LCA + \(\dist\), binary lifting & 16 & \lv{4} & ``the inner loop the compression calls'' & \chref{c11}, \thmref{c11-distf} \\
Euler \(\tin/\tout\), ancestor test & 5 & \lv{3} & ``the sort key and the stack predicate'' & \chref{c24} \\
Bottom-up subtree fold & 13 & \lv{4} & ``run the same DP on fewer nodes'' & \chref{c10} \\
Weighted centroid / 1-median & 24 & \lv{5} & ``minimise the sum of distances to a set'' & \chref{c18} \\
\midrule
\multicolumn{5}{l}{\itshape Chapter 26 -- new this chapter} \\
Virtual tree \(A(S)\) & 26.1 & \lv{4} & ``compress a fixed tree to a marked set'' & \defref{c30-vt}, \coderef{c30-build} \\
Adjacent-pair closure, \(|A(S)|\le2k-1\) & 26.1 & \lv{4} & ``\(\bigO{k}\) skeleton, independent of \(n\)'' & \thmref{c30-size} \\
Weighted fold on \(A(S)\) & 26.2 & \lv{5} & ``a \(\bigO{k}\) DP that equals the \(\bigO{n}\) one'' & \thmref{c30-fold}, \coderef{c30-fold} \\
Pairwise distance sum & 26.3 & \lv{5} & ``sum \(\dist\) over all pairs of a set'' & \thmref{c30-pairsum}, \coderef{c30-pairsum} \\
Steiner-on-a-tree & 26.3 & \lv{5} & ``minimum connected subgraph spanning marked nodes'' & \secref{c30-apps} \\
Recognition gate & 26.4 & \lv{4} & ``is this even a virtual tree?'' & \defref{c30-signature} \\
Method-boundary pick & 26.5 & \lv{5} & ``what does the compression throw away?'' & \defref{c30-pick} \\
Virtual-tree kit (one listing) & 26 & \lv{5} & ``all the skeletons in one place'' & \coderef{c30-kit} \\
Summit: divisor-tree VT + centroid & Summit & \lv{6} & ``infinite tree, factorial marks, assembly node'' & \coderef{c30-summit} \\
\bottomrule
\end{tabular}
\end{center}

\chaptersection{Related Topics}

\begin{itemize}[nosep, leftmargin=1.4em]
  \item \chref{c11} (Ch.~16, LCA) --- backward, \textbf{promise discharged}: Ch.~16 ended by
        promising that a virtual tree compresses a \(k\)-node set via sort-by-tin plus adjacent-pair
        LCAs, with \(\thmref{c11-distf}\) and lifting as the inner loop. \(\secref{c30-construction}\)
        is exactly that build -- it \emph{calls} \(\lca\) \(O(k)\) times and never re-derives it.
  \item \chref{c24} (Ch.~5, Euler Tour / \(\tin/\tout\)) --- backward: the DFS order is the sort key
        and the ancestor predicate of the stack; the flatten is reused, not re-taught.
  \item \chref{c10} (Ch.~13, Tree DP) --- backward: the bottom-up subtree fold is run verbatim on the
        compressed tree (\(\thmref{c30-fold}\)); 26.2 adds only the edge-weight and \(\mathit{cnt}\) carry.
  \item \chref{c23} (Ch.~25, Link-Cut Trees) --- contrast and boundary: link-cut \emph{moves} a forest,
        a virtual tree \emph{shrinks} a fixed one; 26.5 hands every ``edges change'' spec to \chref{c23}
        (\(\defref{c30-pick}\)). Ch.~25's mixed set forward-cited CF 809E to this chapter.
  \item \chref{c18} (Ch.~24, Centroid Decomposition) --- forward/tool: the Summit's assembly node is a
        \emph{weighted centroid} (1-median), the same minimiser Ch.~24 uses for distance sums, now found
        on the virtual tree rather than the full one.
  \item \chref{c16} (Ch.~22, HLD) / \chref{c24} (Ch.~5) --- contrast: path/subtree \emph{ranges} on a
        fixed tree do not shrink with \(k\), so they stay with HLD/Euler; a virtual tree is for a
        marked set, not an interval (\(\secref{c30-boundary}\)).
  \item \chref{c21} (Ch.~27, Mo's Algorithm) --- forward: an offline batch of overlapping node-set
        queries with no inverse is Mo-on-trees, complementing the per-set virtual-tree fold
        (\(\secref{c30-boundary}\)). \textbf{This is the Next chapter.}
  \item Bitmask Steiner DP (outside Part~V) --- boundary: connecting terminals that may meet at new
        nodes is Dreyfus--Wagner \(O(3^k n)\), not a virtual tree; mixed CF 152E is the honest contrast.
\end{itemize}

\begin{lastword}[Last Word --- Chapter 26]
\emph{``A virtual tree is a fixed tree seen through a keyhole: sort a marked set by \(\tin\), add the
LCAs of adjacent pairs, and a monotone stack hands you a \(\bigO{k}\) skeleton whose compressed edges
remember their real lengths -- so every distance, count and connected structure among the marked nodes
survives the shrinking, and an \(\bigO{n}\) DP becomes \(\bigO{k\log k}\).''}

\medskip
Chapter~16 promised this compression and Chapter~26 paid it. Section 26.1 closed the marked set under
\(\lca\) and bounded it at \(2k-1\) (\(\thmref{c30-size}\)); 26.2 ran Chapter~13's fold on the shrunken
tree with the two carrying rules that make it faithful (\(\thmref{c30-fold}\)); 26.3 cashed it in on
the two workhorse families -- the pairwise sum \(\sum w\cdot\mathit{cnt}(k-\mathit{cnt})\) and the
forced tree-Steiner tree (\(\thmref{c30-pairsum}\)); 26.4 drilled the recognition gate so you reach for
the tool only when it fits (\(\defref{c30-signature}\)); and 26.5 drew the honest boundary where the
compression loses information and hands the work to tree DP, bitmask Steiner, HLD or link-cut
(\(\defref{c30-pick}\)). The Summit CF 1292D (\(\probref{c30-summit}\)) took the idea to its limit -- an
infinite tree you compress to a decreasing-prime-factor trie and solve with a weighted centroid.

\medskip
Carry three reflexes forward:
\begin{enumerate}[nosep, leftmargin=1.6em]
  \item \emph{Sort by \(\tin\), then close.} The \(k-1\) adjacent-pair LCAs are all you need
        (\(\thmref{c30-size}\)); never pay \(\binom{k}{2}\).
  \item \emph{Weight the edges by real length.} A compressed edge is worth \(\depth(v)-\depth(u)\);
        a unit weight is the quiet WA that breaks every distance downstream
        (\(\defref{c30-edgeweight}\)).
  \item \emph{Check the marked set is doing work.} If the answer needs every node, or the terminals
        sit on a graph and meet at new nodes, or the tree is being edited -- you do not have a virtual
        tree (\(\defref{c30-pick}\)).
\end{enumerate}

\medskip
And three habits: dedup after inserting LCAs so the stack invariant holds; root the virtual tree at the
global LCA, not a marked node; and price the run by \(\sum k\) versus \(n\) \emph{before} coding, so you
only build \(A(S)\) when it actually shrinks. You now own the one tree engine that answers by making a
big tree \emph{small} -- the static, compressive companion to every moving or path-based structure in Part~V.

\medskip
\textbf{Next:} \chref{c21} (Mo's Algorithm on Trees) turns the difficulty inside-out again. Where a
virtual tree re-derives the answer per marked set in \(\bigO{k\log k}\), Mo orders a large batch of
\emph{offline} queries so that consecutive sets differ by one node, amortising the cost to a single
add/remove step each -- the technique for when the queries overlap so heavily that rebuilding \(A(S)\)
every time wastes work Mo would share.
\end{lastword}


\closechapter
